\documentclass[a4paper,9pt,twocolumn]{extarticle}
\usepackage[UTF8]{ctex}
\usepackage[top=1cm,bottom=1cm,left=1cm,right=1cm,columnsep=0.6cm]{geometry}
\usepackage{fancyhdr,titlesec,xcolor,amsmath,amssymb,listings}
\usepackage[colorlinks=true,linkcolor=black,urlcolor=black,citecolor=black]{hyperref}
\setmainfont{Times New Roman}\setsansfont{Arial}\setmonofont{Consolas}[Scale=0.9]
\definecolor{codebg}{gray}{0.95}\definecolor{commentgreen}{gray}{0.45}
\lstset{language=C++,basicstyle=\fontsize{7.5}{9}\selectfont\ttfamily\color{black},
  keywordstyle=\bfseries\color{black},commentstyle=\itshape\color{commentgreen},
  stringstyle=\color{black},numbers=left,numberstyle=\tiny\color{gray},
  stepnumber=1,numbersep=4pt,backgroundcolor=\color{codebg},
  showspaces=false,showstringspaces=false,showtabs=false,tabsize=2,
  breaklines=true,breakatwhitespace=false,frame=single,framerule=0.3pt,
  rulecolor=\color{black!30},framesep=4pt,aboveskip=4pt,belowskip=4pt,
  lineskip=0.5pt,xleftmargin=4pt,xrightmargin=4pt}
\pagestyle{fancy}\fancyhf{}\fancyhead[L]{\small\sffamily\color{gray}\leftmark}
\fancyhead[R]{\small\sffamily\color{gray}\rightmark}\fancyfoot[C]{\small\sffamily\thepage}
\renewcommand{\headrulewidth}{0.3pt}\setlength{\headheight}{14pt}
\titleformat{\section}{\Large\bfseries\sffamily}{\thesection}{0.8em}{}[{\vspace{-0.3em}\titlerule[0.3pt]}]
\titleformat{\subsection}{\normalsize\bfseries\sffamily}{\thesubsection}{0.6em}{}
\titlespacing{\section}{0pt}{8pt}{4pt}\titlespacing{\subsection}{0pt}{6pt}{2pt}
\newcommand{\algo}[1]{\subsection{#1}}
\begin{document}
\title{\Huge\bfseries\sffamily ACM/ICPC}\author{\large The Three Stooges}\date{\today}\maketitle
\tableofcontents\newpage

\section{不要忘记}

\algo{土制cph(仅linux)}
\begin{lstlisting}
from os import system as e, listdir as l
from sys import argv
_, q, f = argv  # q: 题号, f: 文件名

if "py" in f:
	r = f"python3 {f}"
else:
	e(f"g++ -std=c++23 -o2 -Wall {f}.cpp -o {f}")
	r = "./" + f

#
d = "samples-" + q.capitalize()
for i in l(d):
	if not "in" in i: continue
	print(i)
	p = d + "/" + i[:-2]
	if e(f"timeout 2 {r}<{p}in>{p}out"): print("TLE or RE")
	else:
		with open(f"{p}out") as o, open(f"{p}ans") as a:
			print("AC" if o.read().split()==a.read().split() else "WA")
\end{lstlisting}

\algo{土制cph(全平台)}
\begin{lstlisting}
from os import listdir as l, system as e
from os.path import join as jp
from sys import argv as a
import subprocess as s
_,q,f = a
if "py" in f:
    cmd = ["python3", f]
else:
    e(f"g++ -std=gnu++20 -O2 -Wall {f}.cpp -o {f}")
    cmd = jp(".", f"{f}.exe")

d='samples-'+q.capitalize()
g,r,p,n='\033[32m','\033[31m','\033[35m','\033[0m'
for i in l(d):
    if not 'in' in i: continue
    t=jp(d,i[:-2])
    print(i,'测评结果 ',end='',flush=1)
    try:
        k = s.run(cmd, timeout=2, stdin=open(f"{t}in"), 
                 stdout=open(f"{t}out", 'w'), stderr=s.PIPE, text=1)
        c = lambda o,a : o.read().split() == a.read().split()
        # c = lambda o,a : all(abs(float(x)-float(y))/max(1,abs(float(y)))<=1e-4 
        #               for x,y in zip(o.read().split(),a.read().split()))
        print(f'{g}AC{n}' if c(open(f"{t}out"),open(f"{t}ans")) else f'{r}WA{n}')
        if k.stderr:print(f'{r}{k.stderr}{n}')
    except s.TimeoutExpired as k:
        print(f'{p}TLE or RE{n}')
        if k.stderr: print(f'{r}{k.stderr}{n}')
\end{lstlisting}

\algo{随机数}
\begin{lstlisting}
mt19937_64 rng(chrono::steady_clock::now().time_since_epoch().count());
int Rand(int x, int y) {
	return uniform_int_distribution<int>(x, y)(rng);
}
shuffle(a.begin(), a.end(), rng);
\end{lstlisting}

\algo{对拍}
\begin{lstlisting}
// checker.cpp
while (1) {
	system("A_gen > _data.in");  // linux 下换成 ./xxx
	system("A_bf < _data.in > A.ans");
	system("A < _data.in > A.ans");
	if (system("fc A.out A.ans")) {  // linux 下换成 diff
		cout << "WA\n"; return 0;
	} else {
		cout << "AC\n";
	}
}
\end{lstlisting}

\algo{解除python大数限制}
\begin{lstlisting}
import sys
sys.set_int_max_str_digits(100005)
\end{lstlisting}

\algo{常用函数}
\begin{lstlisting}
// 快读
inline int read()
{
	int x=0,f=1;char ch=getchar();
	while (ch<'0'||ch>'9'){if (ch=='-') f=-1;ch=getchar();}
	while (ch>='0'&&ch<='9'){x=x*10+ch-48;ch=getchar();}
	return x*f;
}

// 取模
const int MOD=998244353;
ll add(ll x, ll y) { return (x+y)%MOD; }
ll del(ll x, ll y) { return add(x, MOD-y); }
ll mul(ll x, ll y) { return (x*y)%MOD; }
ll qpow(ll a, ll b = MOD - 2) {
	ll res = 1;
	for (; b; b >>= 1) {
		if (b & 1) { res = mul(res, a); }
		a = mul(a, a);
	}
	return res;
}

// 计算一个整数的二进制表示中有多少个 1
__builtin_popcountll(i);
std::popcount((unsigned long long)i);  // c++20

// 计算一个整数的二进制表示中最高位的 1
__lg(i);  // i 最多 32 位
\end{lstlisting}

\algo{高精度}
\begin{lstlisting}
istream& operator>> (istream& is, __int128 &x)
{
	string s; is>>s;
	bool f=false;
	__int128 ans = 0;
	for (auto c: s) {
		if (c == '-') { f = true; }
		if (isdigit(c)) { ans = ans * 10 + (c-'0'); }
	}
	x = f? -ans: ans;
	return is;
}
ostream& operator<< (ostream& os, __int128 &x)
{
	string ans;
	function<void(__int128)> write = [&](__int128 x){
		if (x < 0) { ans += '-'; x=-x; }
		if (x > 9) { write(x/10); }
		ans += '0'+x%10;
	};
	write(x); return os << ans;
}
\end{lstlisting}

\algo{vector相关}
\begin{lstlisting}
// 多维 vector，需要 c++17，c++14 只能老老实实写 vector<vector<int>>
vector a(n+1, vector (n+1, vector<int>(n+1)));

// 最值
int mx = *min_element(vec.begin(), vec.end(), cmp);
int mn = *min_element(vec.begin(), vec.end(), cmp);

// 去重与离散化
for (int i = 1; i <= n; ++i) { vec.push_back(a[i]); }
sort(vec.begin(), vec.end());
vec.erase(unique(vec.begin(), vec.end()), vec.end());
for (int i = 1; i <= n; ++i) {
	idx[i]=lower_bound(vec.begin(), vec.end(), a[i])-vec.begin()+1;  // 下标从1开始
}

// 前缀和
vector <ll> a(n), s(n);
partial_sum(a.begin(), a.end(), s.begin());
ll sum = accumulate(a.begin(), a.end(), 0LL);

// 填充 1~n，令 a[i]=i
iota(a.begin(), a.end(), 0);
\end{lstlisting}

\algo{整数域三分}
\begin{lstlisting}
// 整数域三分求单峰函数最大值
int l, r;
while (l + 2 <= r) {
	int w = (r - l) / 3;
	int m1 = l + w; int v1 = f(m1);
	int m2 = r - w; int v2 = f(m2);
	if (v1 <= v2) {
		l = m1 + 1;
	} else {
		r = m2 - 1;
	}
} 
cout << max(f(l), f(r)) << "\n";
\end{lstlisting}

\algo{运行时间}
\begin{lstlisting}
auto start = chrono::high_resolution_clock::now();
auto stop = chrono::high_resolution_clock::now();
auto duration = std::chrono::duration_cast<std::chrono::milliseconds>(stop - start).count();
if (duration > 998) { cout << ans << "\n"; exit(0); }
\end{lstlisting}

\algo{记得快读快输}
\begin{lstlisting}
ios::sync_with_stdio(0); 
cin.tie(0), cout.tie(0);
\end{lstlisting}

\algo{！！！数据范围}
\begin{lstlisting}
注意爆int

    long long: 8字节，-9,223,372,036,854,775,808到9,223,372,036,854,775,807
    unsigned long long: 8字节，0到18,446,744,073,709,551,615
    如果爆long long 就用__int128,但是注意输入输出

    尽量不要使用double输入，容易超时
\end{lstlisting}

\algo{\_\_int128的使用}
\begin{lstlisting}
#define int __int128
inline void read(int &n){
    int x=0,f=1;
    char ch=getchar();
    while(ch<'0'||ch>'9'){
        if(ch=='-') f=-1;
        ch=getchar();
    }
    while(ch>='0'&&ch<='9'){
        x=(x<<1)+(x<<3)+(ch^48);
        ch=getchar();
    }
    n=x*f;
}
inline void print(int n){
    if(n<0){
        putchar('-');
        n*=-1;
    }
    if(n>9) print(n/10);
    putchar(n % 10 + '0');
}
#undef int
\end{lstlisting}

\algo{快读模板}
\begin{lstlisting}
#define re register
#define il inline
il int read()
{
    re int x=0,f=1;char c=getchar();
    while(c<'0'||c>'9'){if(c=='-') f=-1;c=getchar();}
    while(c>='0'&&c<='9') x=(x<<3)+(x<<1)+(c^48),c=getchar();
    return x*f;
}
\end{lstlisting}

\algo{随机数生成器}
\begin{lstlisting}
mt19937_64 rng(chrono::steady_clock::now().time_since_epoch().count());
int Rand(int x, int y) {
	return uniform_int_distribution<int>(x, y)(rng);
}
shuffle(a.begin(), a.end(), rng);
\end{lstlisting}

\algo{随机图}
\begin{lstlisting}
// n 为节点个数
// m 为边的个数，默认为 -1，表示建树
// root 为根节点编号，默认为 -1，
vector<array<int,2>> Graph(int n, int m=-1, int root=-1)
{
	assert(m==-1 || (n-1<=m && m<=1LL*n*(n-1)/2));
	assert(root==-1 || (1<=root && root<=n));
	vector<array<int,2>> E;
	set<array<int,2>> diff;
	// 先建一棵以 0 为根的树，然后选根偏移
	for (int i = 1; i < n; ++i) {
		E.push_back( { i, Rand(0, i-1) } );
	}
	if (root == -1) { root = Rand(0, n-1); }
	for (auto& [x, y]: E) {
		x = (x+root) % n + 1;
		y = (y+root) % n + 1;
		diff.insert( { x, y } );
		diff.insert( { y, x } );  // 无向图中需要添加这条反向边
	}
	// 从树建图
	if (m != -1) {
		for (int i = n-1; i < m; ++i) {
			while (1) {
				int x = Rand(1, n);
				int y = Rand(1, n);
				if (x == y || diff.count({x,y})) { continue; }
				E.push_back( { x, y } );
				diff.insert( { x, y } );
				diff.insert( { y, x } );  // 无向图中需要添加这条反向边
				break;
			}
		}
	}
	shuffle(E.begin(), E.end(), rng);
	return E;
}
\end{lstlisting}

\algo{对拍器}
\begin{lstlisting}
#include<bits/stdc++.h>
using namespace std;
int main(){
    int t=0;
    while(true){
    cout<<"test: "<<t++<<endl;
    system("data.exe>data.in");//data.exe 随机数生成器
    system("std.exe<data.in>data.out");//std.exe 标准程序
    system("solve.exe<data.in>solve.out");//solve.exe 暴力求解算法
    if(system("fc std.out solve.out>diff.log")){
        cout<<"WA"<<endl;
        break;
    } 
}
cout<<"AC\n";
}
/*
// checker.cpp
while (1) {
	system("A_gen > _data.in");  // linux 下换成 ./xxx
	system("A_bf < _data.in > A.ans");
	system("A < _data.in > A.ans");
	if (system("fc A.out A.ans")) {  // linux 下换成 diff
		cout << "WA\n"; return 0;
	} else {
		cout << "AC\n";
	}
}*/
\end{lstlisting}

\section{数据结构}

\algo{对顶multiset}
\begin{lstlisting}
struct PairingMultiset {
    int sz;
    multiset <int> minH, maxH;
    ll minS, maxS;
    void clear()
    {
        minH.clear(); maxH.clear();
        sz = minS = maxS = 0;
    }
    void move()
    {
        while ((int)minH.size() < (sz+1)/2) {
            minS += *maxH.rbegin();
            minH.insert(*maxH.rbegin());
            maxS -= *maxH.rbegin();
            maxH.erase(maxH.find(*maxH.rbegin()));
        }
        while ((int)minH.size() > (sz+1)/2) {
            maxS += *minH.begin();
            maxH.insert(*minH.begin());
            minS -= *minH.begin();
            minH.erase(minH.find(*minH.begin()));
        }
    }
    void add(int x)
    {  // 向对顶堆中插入一个元素
        ++sz;
        if (minH.empty() || x >= *minH.begin()) {
            minS += x;
            minH.insert(x);
        } else {
            maxS += x;
            maxH.insert(x);
        }
        move();
    }
    void del(int x)
    {  // 从对顶堆中删除一个x
        --sz;
        if (x >= *minH.begin()) {
            minS -= x;
            minH.erase(minH.find(x));
        } else {
            maxS -= x;
            maxH.erase(maxH.find(x));
        }
        move();
    }
    ll med()
    {  // 求中位数
        return *minH.begin();
    }
    ll dis_to_med()
    {  // 求对顶堆中的数字到中位数的距离和
        ll minSZ=(sz+1)/2, maxSZ=sz-minSZ, m=med();
        return (minS-minSZ*m)+(maxSZ*m-maxS);
    }
} p;
\end{lstlisting}

\algo{带权并查集}
\begin{lstlisting}
struct DSU {
    int n;
    vector<int> f, siz, pre;
    DSU(int n): n(n),
    			f(vector<int>(n+1)),
    			siz(vector<int>(n+1,1)),
    			pre(vector<int>(n+1))
    {
        iota(f.begin(), f.end(), 0);
    }
    int find(int x) {
        if (f[x] == x) { return x; }
        int fa = find(f[x]);
        pre[x] += pre[f[x]];
        return f[x]=fa;
    }
    bool same(int x, int y) {
        return find(x) == find(y);
    }
    void merge(int x, int y) {
        x = find(x);
        y = find(y);
        f[x] = y;
        pre[x] += siz[y];
        siz[y] += siz[x];
        siz[x] = 0;
    }
    int query(int x, int y) {
        return abs(pre[x]-pre[y])-1;
    }
};
\end{lstlisting}

\algo{Trie}
\begin{lstlisting}
const string VAL="abcdefghijklmnopqrstuvwxyzABCDEFGHIJKLMNOPQRSTUVWXYZ0123456789";
const int S=3e6+5, V=62;

struct Trie {
    unordered_map<char,int> p;
    int tot, ch[S][V], pass[S];
    int newNode() {
        ++tot;
        memset(ch[tot], 0, sizeof(ch[tot]));
        pass[tot] = 0;
        return tot;
    }
    void init() {
        tot = -1;
        newNode();
    }
    Trie() {
        for (int i = 0; i < V; ++i) { p[VAL[i]] = i; }
        init();
    }
    void insert(string s) {
        int now = 0;
        for (auto c: s) {
            if (!ch[now][p[c]]) { ch[now][p[c]] = newNode(); }
            now = ch[now][p[c]];
            ++pass[now];
        }
    }
    int query(string s) {
        int now = 0;
        for (auto c: s) {
            if (!ch[now][p[c]]) { return 0; }
            now = ch[now][p[c]];
        }
        return pass[now];
    }
} trie;
\end{lstlisting}

\algo{ST表}
\begin{lstlisting}
template <typename T>
struct ST {
    int n=0, I=0;
    vector<int> Log;
    vector<vector<T>> st;
    ST() { }
    ST(const vector<T>& a): n(a.size()), Log(vector<int>(n+1)) {
        for (int i = 2; i <= n; ++i) {
        	I=Log[i]=Log[i/2]+1;
        }
        st.assign(I+1, vector<T>(n));
        copy(a.begin(), a.end(), st[0].begin());
        for (int i = 1; i <= I; ++i) {
            for (int j = 0; j+(1<<(i-1)) < n; ++j) {
                st[i][j] = max( st[i-1][j], st[i-1][j+(1<<(i-1))] );
            }
        }
    }
    T query(int l, int r) {
    	assert(l <= r);  // l > r 时的返回值需要特殊定义
        int s = Log[r-l+1];
        return max( st[s][l], st[s][r-(1<<s)+1] );
    }
    int find(int l, T x) {  // 第一个区间 [l,r] 最值大于等于 x 的 r
    	if (l >= n) { return -1; }
        int rl=l-1, rr=n;
        while (rl != rr-1) {
            int mid = (rl + rr) >> 1;
            if (query(l, mid) >= x) { rr = mid; }
            else { rl = mid; }
        }
        if (rr < n) { return rr; }
        return -1;
    }
	int find(int l, int x) {  // 倍增写法
	    assert(l < n);
	    if (st[0][l] < x) { return l; }
	    if (query(l, n-1) >= x) { return n; }
	    int ans = l, g = st[0][l];
	    for (int i = I; i >= 0; --i) {
	        if (ans+(1<<i) < n && gcd(g, st[i][ans]) >= x) {
	            g = max(g, st[i][ans]);
	            ans += 1 << i;
	        }
	    }
	    return ans;
	}
};
\end{lstlisting}

\algo{树状数组}
\begin{lstlisting}
struct BIT {
	int n;
    vector<int> t;
    BIT(int n): n(n), t(n+1) { }  // 注意值域树状数组中 n=tot
    int lowbit(int x) { return x&-x; }
    void modify(int x, int d) {
        for (; x <= n; x += lowbit(x)) { t[x] += d; }
    }
	// 前缀和
    int query(int x) {
        int res = 0;
        for (; x; x -= lowbit(x)) { res += t[x]; }
        return res;
    }
    // 求最小的前缀和等于x的位置
	int select(int x) {
	    int ans = 0, sum = 0;
	    for (int i = __lg(n); i >= 0; --i) {
	        if (ans+(1<<i) <= n && sum+t[ans+(1<<i)] < x) { 
	            ans += 1<<i;
	            sum += t[ans];
	        }
	    }
	    return ans + 1;
	}
};
BIT T(n);
\end{lstlisting}

\algo{线段树}
\begin{lstlisting}
#define lson (p << 1)
#define rson (p << 1 | 1)
#define m ((l+r)>>1)
template<typename Info>
struct segmentTree {
    int n;
    vector<Info> a;
    
    void push_up(int p) {
        a[p] = a[lson] + a[rson];
    }

    template<typename F>
    Info dfs(int p, int l, int r, int x, int y, F&& op) {
        if (y <= l || r <= x) { return Info(); }
        if (x <= l && r <= y) { op(p,l,r); return a[p]; }
        Info res;
        res = res + dfs(lson, l, m, x, y, op);
        res = res + dfs(rson, m, r, x, y, op);
        push_up(p);
        return res;
    }

	template<typename F>
	pair<int,Info> findFirst(int p, int l, int r, int x, int y, F&& pred) {
	    if (y <= l || r <= x) { return {-1,Info()}; }
	    if (x <= l && r <= y && !pred(a[p])) { return {-1,Info()}; }
	    if (l == r-1) { return {l,a[p]}; }
	    pair<int,Info> res = findFirst(lson, l, m, x, y, pred);
	    if (res.first == -1) {
	        res = findFirst(rson, m, r, x, y, pred);
	    }
	    return res;
	}
	
	template<typename F>
	pair<int,Info> findLast(int p, int l, int r, int x, int y, F&& pred) {
	    if (y <= l || r <= x) { return {-1,Info()}; }
	    if (x <= l && r <= y && !pred(a[p])) { return {-1,Info()}; }
	    if (l == r-1) { return {l,a[p]}; }
	    pair<int,Info> res = findLast(lson, m, r, x, y, pred);
	    if (res.first == -1) {
	        res = findLast(rson, l, m, x, y, pred);
	    }
	    return res;
	}

    segmentTree(int n = 0): n(n), a(4 << __lg(n)) { }
    segmentTree(vector<Info>&& b): segmentTree(b.size()) {
        function<void(int,int,int)> build = [&](int p, int l, int r) -> void {
            if (r - l == 1) { a[p] = Info(b[l]); return; }
            build(lson, l, m);
            build(rson, m, r);
            push_up(p);
        };
        build(1, 0, n);
    }

    void modify(int pos, int x) {
        dfs(1, 0, n, pos, pos+1, [&](int p, int l, int r){ a[p].mx = x; });
    }

    int query(int l, int r) {
        return dfs(1, 0, n, l, r+1, [](int,int,int){}).mx;
    }

    void Debug(int p, int l, int r) {
        cerr << format("a[{}] = [ {}, {}, {} ]\n", p, l, r, a[p].mx);
        if (l == r - 1) { return; }
        Debug(lson, l, m);
        Debug(rson, m, r);
    }
    void Debug() { Debug(1, 0, n); }
};
#undef lson
#undef rson
#undef m
struct Info {
    int mx;
    Info(int mx = 0): mx(mx) { }
    friend Info operator+ (Info u, Info v) {
        return Info( max(u.mx, v.mx) );
    }
};
\end{lstlisting}

\algo{逆序对}
\begin{lstlisting}
ll msort(vector<int> &a, int l, int r)
{
    if (l == r) { return 0; }
    int mid = (l + r) >> 1;
    ll ans=0;
    ans += msort(a, l, mid);
    ans += msort(a, mid+1, r);
    int i=l, j=mid+1;
    vector<int> b;
    while (i <= mid && j <= r) {
        if (a[i] <= a[j]) { b.push_back(a[i++]); }
        else              { b.push_back(a[j++]); ans+=mid-i+1; }
    }
    while (i <= mid) { b.push_back(a[i++]); }
    while (j <= r)   { b.push_back(a[j++]); }
    copy(b.begin(), b.end(), a.begin()+l);
    return ans;
}
\end{lstlisting}

\algo{笛卡尔树}
\begin{lstlisting}
vector<int> stk;
for (int i = 1; i <= n; ++i) {
	int flag = 0;
	while (!stk.empty()) {
		if (a[stk.back()] > a[i]) {  // 小根堆
			flag = stk.back(); stk.pop_back();
		} else {
			break;
		}
	}
	if (flag) { L[i] = flag; }
	if (!stk.empty()) { R[stk.back()] = i; }
	stk.push_back(i);
}
function<void(int)> dfs = [&](int x){
	cerr << "now node " << x << "\n";
	cerr << ":: L = " << L[x] << "\n";
	cerr << ":: R = " << R[x] << "\n";
	if (L[x]) dfs(L[x]);
	if (R[x]) dfs(R[x]);
};
dfs(stk[0]);
\end{lstlisting}

\algo{可持久化}
\begin{lstlisting}
5 10
59 46 14 87 41
0 2 1
0 1 1 14
0 1 1 57
0 1 1 88
4 2 4
0 2 5
0 2 4
4 2 1
2 2 2
1 1 5 91
---
59
87
41
87
88
46
\end{lstlisting}

\algo{线段树模板}
\begin{lstlisting}
//区间乘法
#include<bits/stdc++.h>
using namespace std;
#define maxn 1000010
#define ll long long
int p;
ll a[maxn], w[maxn*4];
ll lzy_add[maxn*4], lzy_mul[maxn*4];

void pushup(const int u) {
    w[u] = (w[u<<1] + w[u<<1|1]) % p;
}

void build(const int u, int L, int R) {
    lzy_mul[u] = 1;
    lzy_add[u] = 0;
    if (L == R) {
        w[u] = a[L] % p; 
        return;
    }
    int M = (L+R)>>1;
    build(u<<1, L, M);
    build(u<<1|1, M+1, R);
    pushup(u);
}

bool InRange(int L, int R, int l, int r) {
    return (l <= L) && (R <= r);
}

bool OutofRange(int L, int R, int l, int r) {
    return (R < l) || (L > r);
}

void maketag(int u, int L, int R, ll x, int type) {
    if (type == 1) {
        (lzy_mul[u] *= x) %= p;
        (lzy_add[u] *= x) %= p;
        (w[u] *= x) %= p;
    } else {
        (lzy_add[u] += x) %= p; 
        (w[u] += (R-L+1) * x) %= p;
    }
}

void pushdown(int u, int L, int R) {
    int M = (L+R)>>1;
    // 左子节点
    maketag(u<<1, L, M, lzy_mul[u], 1);
    maketag(u<<1, L, M, lzy_add[u], 2);
    // 右子节点
    maketag(u<<1|1, M+1, R, lzy_mul[u], 1);
    maketag(u<<1|1, M+1, R, lzy_add[u], 2);
    lzy_mul[u] = 1;
    lzy_add[u] = 0;
}

ll query(int u, int L, int R, int l, int r) {
    if (InRange(L, R, l, r)) {
        return w[u];
    } else if (!OutofRange(L, R, l, r)) {
        int M = (L+R)>>1;
        pushdown(u, L, R);
        return (query(u<<1, L, M, l, r) + query(u<<1|1, M+1, R, l, r)) % p; 
    } else {
        return 0;
    }
}

void update(int u, int L, int R, int l, int r, ll x, int type) {
    if (InRange(L, R, l, r)) {
        maketag(u, L, R, x, type);
    } else if (!OutofRange(L, R, l, r)) {
        int M = (L+R)>>1;
        pushdown(u, L, R);
        update(u<<1, L, M, l, r, x, type);
        update(u<<1|1, M+1, R, l, r, x, type);
        pushup(u);
    }
}

int main() {
    ios::sync_with_stdio(false);
    cin.tie(nullptr);
    ll n, m;
    cin >> n >> m >> p;
    for (int i = 1; i <= n; i++) 
        cin >> a[i];
    build(1, 1, n);
    while (m--) {
        int op, x, y;
        cin >> op >> x >> y;
        if (op == 1 || op == 2) {
            ll k;
            cin >> k;
            update(1, 1, n, x, y, k, op); 
        } else {
            cout << query(1, 1, n, x, y) << endl; 
        }
    }
    return 0;
}
\end{lstlisting}

\algo{主席树（可持久化线段树）}
\begin{lstlisting}
//对于单点历史修改与查询
#include<bits/stdc++.h>
using namespace std;
#define int long long
const int maxn=(1e6+1);
int top,root[maxn],arr[maxn],n,m,cnt=0;
struct kkk{
    int l,r,val;
} a[maxn*30];
int build(int l,int r){
    int rt=++cnt;
    if(l==r)a[rt].val=arr[l];
    else {
        int mid=(l+r)>>1;
        a[rt].l=build(l,mid);
        a[rt].r=build(mid+1,r);
    }
    return rt;
}
int update(int ji,int jv,int l,int r,int i)//单点修改
{
    int rt=++cnt;
    a[rt].l=a[i].l;
    a[rt].r=a[i].r;
    if(l==r) a[rt].val=jv;
    else {
        int mid=(l+r)>>1;
        if(ji<=mid)a[rt].l=update(ji,jv,l,mid,a[rt].l);
        else a[rt].r=update(ji,jv,mid+1,r,a[rt].r);
    }
    return rt;
}
int query(int ji,int l,int r,int i){
    if(l==r) return a[i].val;
    int mid=(l+r)>>1;
    if(ji<=mid) return query(ji,l,mid,a[i].l);
    else return query(ji,mid+1,r,a[i].r);
}
signed main(){
    ios_base::sync_with_stdio(false);
    cin.tie(nullptr);
    cin>>n>>m;
    for(int i=1;i<=n;i++){
        cin>>arr[i];
    }
    root[0]=build(1,n);
    for(int i=1,version,op,x,v;i<=m;i++){
        cin>>version>>op>>x;
        if(op==1){
            cin>>v;
            root[i]=update(x,v,1,n,root[version]);
        }else {
            root[i]=root[version];
            cout<<query(x,1,n,root[i])<<endl;
        }
    }
}
\end{lstlisting}

\algo{开点线段树}
\begin{lstlisting}
//适用于范围很大，但是操作较少
//用于减小空间复杂度
//如果操作过多，则与没有优化相同，因此可能被卡
//在线段树的合并与分裂中不可替代，树链剖分
#include<bits/stdc++.h>
using namespace std;
typedef long long ll;
const int maxn=100001;
ll le[maxn],ri[maxn],sum[maxn],add[maxn],cnt;
void up(ll i, ll l, ll r){
    sum[i]=sum[l]+sum[r];
}
void lazy(ll i,ll v, ll n){
    sum[i]+=v*n;
    add[i]+=v;
}
void down(ll i, ll ln, ll rn){
    if(add[i]!=0){
     if(le[i]==0){
        le[i]=++cnt;
     }
    if(ri[i]==0){
        ri[i]=++cnt;
    }
    lazy(le[i],add[i],ln);
    lazy(ri[i],add[i],rn);
    add[i]=0;
    }
}
void Add(ll jl,ll jr,ll v,ll l,ll r,ll i){
     if(jl<=l&&r<=jr){
        lazy(i,v,r-l+1);
       return;
     }
     
    ll mid=(l+r)>>1;
     down(i,mid-l+1,r-mid);
     if(jl<=mid){
        if(le[i]==0){
            le[i]=++cnt;
        }
        Add(jl,jr,v,l,mid,le[i]);
     }
     if(jr>mid){
        if(ri[i]==0){
            ri[i]=++cnt;
        }
        Add(jl,jr,v,mid+1,r,ri[i]);
     }
    
    up(i,le[i],ri[i]);
}
ll query(ll jl,ll jr,ll l, ll r,ll i){
    if(jl<=l&&r<=jr){
        return sum[i];
    }
    ll mid=(l+r)>>1;
    down(i,mid-l+1,r-mid);
    ll ans=0;
    if(jl<=mid){
        if(le[i]!=0){
            ans+=query(jl,jr,l,mid,le[i]);
        }
    }
    if(jr>mid){
        if(ri[i]!=0){
            ans+=query(jl,jr,mid+1,r,ri[i]);
        }
    }
    return ans;
}
int main(){
    int n,m;
    cin>>n>>m;
    cnt=1;
    while(m--){
        int opt;
        cin>>opt;
        if(opt==1){
            int l,r,v;
            cin>>l>>r>>v;
            Add(l,r,v,1,n,1);
        }
    else {
        int l,r;
        cin>>l>>r;
        cout<<query(l,r,1,n,1)<<endl;
    }
}
}
\end{lstlisting}

\algo{权值线段树}
\begin{lstlisting}
权值线段树模板
#define lson pos<<1
#define rson pos<<1|1
void build(int pos,int l,int r)
{
	int mid=(l+r)>>1;
	if(l==r)
	{
		tree[pos]=a[l];//a[l]表示数l有多少个
		return;
	}
	build(lson,l,mid);
	build(rson,mid+1,r);
	tree[pos]=tree[lson]+tree[rson];
}
void update(int pos,int l,int r,int k,int cnt)//表示数k的个数多cnt个
{
	int mid=(l+r)>>1;
	if(l==r)
	{
		tree[pos]+=cnt;
		return;
	}
	if(k<=mid)
		update(lson,l,mid,k,cnt);
	else	
		update(rson,mid+1,r,k,cnt);
	tree[pos]=tree[lson]+tree[rson];
}
int query(int pos,int l,int r,int k)//查询数k有多少个
{
	int mid=(l+r)>>1;
	if(l==r)
		return tree[pos];
	if(k<=mid)
		return query(lson,l,mid,k);
	else
		return query(rson,mid+1,r,k);
}
int kth(int pos,int l,int r,int k)//查询第k大值是多少
{
	int mid=(l+r)>>1;
	if(l==r)
		return l;
	if(k<mid)
		return kth(rson,mid+1,r,k);
	else
		return kth(lson,l,mid,k-mid);
}
\end{lstlisting}

\algo{线段树上最值操作}
\begin{lstlisting}
//基本的实现
#include<bits/stdc++.h>
using namespace std;
const int maxn=1000001;
const int minest=-1000001;
long long arr[maxn<<2],sum[maxn<<2],mx[maxn<<2],sem[maxn<<2],cnt[maxn<<2];
void up(int i){
    int l=i<<1;
    int r=i<<1|1;
    sum[i]=sum[l]+sum[r];
    mx[i]=max(mx[l],mx[r]);
    if(mx[l]>mx[r]){
        cnt[i]=cnt[l];
        sem[i]=max(sem[l],mx[r]);
    }else if(mx[r]>mx[l]){
        cnt[i]=cnt[r];
        sem[i]=max(sem[r],mx[l]);
    }else {
        sem[i]=max(sem[l],sem[r]);
        cnt[i]=cnt[l]+cnt[r];
    }
}
void Lazy(int i,int v){
    if(v<mx[i]){
        sum[i]-=(mx[i]-v)*cnt[i];
        mx[i]=v;
    }
}
void down(int i){
    Lazy(i<<1,mx[i]);
    Lazy(i<<1|1,mx[i]);
}
void setmin(int jl,int jr,int v,int l,int r,int i){
    if(v>mx[i]) return;
    if(jl<=l&&r<=jr&&v>sem[i]){
        Lazy(i,v);
        return;
    }else {
        down(i);
        int mid=(l+r)>>1;
        if(jl<=mid) setmin(jl,jr,v,l,mid,i<<1);
        if(jr>mid) setmin(jl,jr,v,mid+1,r,i<<1|1);
    }
    up(i);
}
long long querymax(int jl,int jr,int l,int r,int i){
    if(jl<=l&&r<=jr){
        return mx[i];
    }
    down(i);
    int mid=(l+r)>>1;
    long long lm=minest,rm=minest;
    if(jl<=mid) lm=querymax(jl,jr,l,mid,i<<1);
    if(jr>mid)  rm=querymax(jl,jr,mid+1,r,i<<1|1);
    return max(lm,rm);
}
long long querysum(int jl,int jr,int l,int r,int i){
    if(jl<=l&&r<=jr){
        return sum[i];
    }
    down(i);
    int mid=(l+r)>>1;
    long long ans=0;
    if(jl<=mid) ans+=querysum(jl,jr,l,mid,i<<1);
    if(jr>mid)  ans+=querysum(jl,jr,mid+1,r,i<<1|1);
    return ans;
}
void build(int l,int r,int i){
    if(l==r){
        sum[i]=arr[l];
        mx[i]=arr[l];
        sem[i]=minest;
        cnt[i]=1;
        return;
    }
    int mid=(l+r)>>1;
    build(l,mid,i<<1),build(mid+1,r,i<<1|1);
    up(i);
}
int main(){
    ios_base::sync_with_stdio(false);
    cin.tie(nullptr);
    int t;
    cin>>t;
    while(t--){
        int n,m;
        cin>>n>>m;
        for(int i=1;i<=n;i++){
            cin>>arr[i];
        }
        build(1,n,1);
        while(m--){
            int opt;
            cin>>opt;
            if(opt==0){
                int l,r;
                long long mi;
                cin>>l>>r>>mi;
                setmin(l,r,mi,1,n,1);
            }else if(opt==1){
                int l,r;
                cin>>l>>r;
                cout<<querymax(l,r,1,n,1)<<endl;
            }else {
                int l,r;
                cin>>l>>r;
                cout<<querysum(l,r,1,n,1)<<endl;
            }
        }
    }
}
\end{lstlisting}

\algo{线段树上序列操作}
\begin{lstlisting}
//线段树上的序列操作，稍难一点
//求连续1的最长子列，有点难度，容易打错

#include<bits/stdc++.h>
using namespace std;
const int maxn=100001;
int arr[maxn],sum[maxn<<2],pre0[maxn<<2],len0[maxn<<2],suf0[maxn<<2],pre1[maxn<<2],len1[maxn<<2],suf1[maxn<<2];
int lazy[maxn<<2];
bool update[maxn<<2];
bool rever[maxn<<2];
void up(int i,int ln,int rn){
    int l=i<<1;
    int r=i<<1|1;
    sum[i]=sum[l]+sum[r];
    len0[i]=max(max(len0[l],len0[r]),suf0[l]+pre0[r]);
    pre0[i]=len0[l]<ln?pre0[l]:pre0[l]+pre0[r];
    suf0[i]=len0[r]<rn?suf0[r]:suf0[l]+suf0[r];
    len1[i]=max(max(len1[l],len1[r]),suf1[l]+pre1[r]);
    pre1[i]=len1[l]<ln?pre1[l]:pre1[l]+pre1[r];
    suf1[i]=len1[r]<rn?suf1[r]:suf1[l]+suf1[r];
}
void updateLazy(int i,int v,int n){
    sum[i]=v*n;
    pre0[i]=suf0[i]=len0[i]=v==0?n:0;
    pre1[i]=suf1[i]=len1[i]=v==1?n:0;
    lazy[i]=v;
    update[i]=true;
    rever[i]=false;
}
void reverseLazy(int i,int n){
    sum[i]=n-sum[i];
    swap(pre0[i],pre1[i]);
    swap(len0[i],len1[i]);
    swap(suf0[i],suf1[i]);
    rever[i]=!rever[i];
}
void down(int i,int ln,int rn){
    if(update[i]){
        updateLazy(i<<1,lazy[i],ln);
        updateLazy(i<<1|1,lazy[i],rn);
        update[i]=false;
    }
    if(rever[i]){
        reverseLazy(i<<1,ln);
        reverseLazy(i<<1|1,rn);
        rever[i]=false;
    }
}
void build(int l,int r,int i){
    if(l==r){
        sum[i]=arr[l];
        pre0[i]=arr[l]?0:1;
        suf0[i]=arr[l]?0:1;
        len0[i]=arr[l]?0:1;
        pre1[i]=arr[l]?1:0;
        suf1[i]=arr[l]?1:0;
        len1[i]=arr[l]?1:0;
        return;
    }
    int mid=(l+r)>>1;
    build(l,mid,i<<1),build(mid+1,r,i<<1|1);
    up(i,mid-l+1,r-mid);
    rever[i]=0;
    update[i]=0;
}
void Update(int jl,int jr,int v,int l,int r,int i){
    if(jl<=l&&r<=jr){
        updateLazy(i,v,r-l+1);
    }else {
        int mid=(l+r)>>1;
        down(i,mid-l+1,r-mid);
        if(jl<=mid) Update(jl,jr,v,l,mid,i<<1);
        if(jr>mid) Update(jl,jr,v,mid+1,r,i<<1|1);
        up(i,mid-l+1,r-mid);
    }
}
void Reverse(int jl,int jr,int l,int r,int i){
    if(jl<=l&&r<=jr){
        reverseLazy(i,r-l+1);
    }else {
        int mid=(l+r)>>1;
        down(i,mid-l+1,r-mid);
        if(jl<=mid) Reverse(jl,jr,l,mid,i<<1);
        if(jr>mid) Reverse(jl,jr,mid+1,r,i<<1|1);
        up(i,mid-l+1,r-mid);
    }
}
tuple<int,int,int> querylength(int jl,int jr,int l,int r,int i){
    if(jl<=l&&r<=jr){
        return make_tuple(len1[i],pre1[i],suf1[i]);
    }
    int mid=(l+r)>>1;
    down(i,mid-l+1,r-mid);
    if(jr<=mid) return querylength(jl,jr,l,mid,i<<1);
    if(jl>mid) return querylength(jl,jr,mid+1,r,i<<1|1);
    tuple<int,int,int> l1=querylength(jl,jr,l,mid,i<<1);
    tuple<int,int,int> r1=querylength(jl,jr,mid+1,r,i<<1|1);
    int llen=get<0>(l1),lpre=get<1>(l1),lsuf=get<2>(l1);
    int rlen=get<0>(r1),rpre=get<1>(r1),rsuf=get<2>(r1);
    int len=max(max(llen,rlen),lsuf+rpre);
    int pre=lpre<mid-max(l,jl)+1?lpre:lpre+rpre;
    int suf=rsuf<min(r,jr)-mid?rsuf:rsuf+lsuf;
    return make_tuple(len,pre,suf);
}
int query(int jl,int jr,int l,int r,int i){
    if(jl<=l&&r<=jr){
        return sum[i];
    }
    int mid=(l+r)>>1;
    down(i,mid-l+1,r-mid);
    int ans=0;
    if(jl<=mid) ans+=query(jl,jr,l,mid,i<<1); 
    if(jr>mid)  ans+=query(jl,jr,mid+1,r,i<<1|1);
    return ans;
}
int main(){
    int n,m;
    cin>>n>>m;
    for(int i=1;i<=n;i++) cin>>arr[i];
    build(1,n,1);
    while(m--){
        int opt,l,r;
        cin>>opt>>l>>r;
        l+=1;
        r+=1;
        if(opt==0){
            Update(l,r,0,1,n,1);
        }else if(opt==1){
            Update(l,r,1,1,n,1);
        }else if(opt==2){
            Reverse(l,r,1,n,1);
        }else if(opt==3){
            cout<<query(l,r,1,n,1)<<endl;
        }else {
            tuple<int,int,int>a=querylength(l,r,1,n,1);
            cout<<get<0>(a)<<endl;
        }
    }
}
\end{lstlisting}

\algo{非经典线段树}
\begin{lstlisting}
// 这类题目一般是可以看出来使用线段树，但是不能满足一般的懒更新的条件
//简单的一般需要进行势能分析，再减枝就可以解决
//例如 洛谷 P4145，此题是开根号，换为取模同理
//P1471 求方差
#include<bits/stdc++.h>
using namespace std;

const int maxn=100001;
double a[maxn],sum[maxn<<2],pf[maxn<<2],lazy[maxn<<2];
void up(int i){
    sum[i]=sum[i<<1]+sum[i<<1|1];
    pf[i]=pf[i<<1]+pf[i<<1|1];
}

void Lazy(int i,double v,int n){
    lazy[i]+=v;
    pf[i]+=2*v*sum[i]+v*v*n;
    sum[i]+=v*n;
}
void down(int i,int ln,int rn){
    if(lazy[i]!=0){
        Lazy(i<<1,lazy[i],ln);
        Lazy(i<<1|1,lazy[i],rn);
        lazy[i]=0;
    }
}
void add(int jl,int jr,double v,int l,int r,int i){
    if(jl<=l&&jr>=r){
        Lazy(i,v,r-l+1);
        return;
    }
    int mid=(l+r)>>1;
    down(i,mid-l+1,r-mid);
    if(jl<=mid){
        add(jl,jr,v,l,mid,i<<1);
    }
    if(jr>mid) add(jl,jr,v,mid+1,r,i<<1|1);
    up(i);
}
double query(int jl,int jr,int jt,int l,int r,int i){
    if(jl<=l&&jr>=r){
        if(jt){
            return pf[i];
        }else return sum[i];
    }
    int mid=(l+r)>>1;
    down(i,mid-l+1,r-mid);
    double ans=0;
    if(jl<=mid) ans+=query(jl,jr,jt,l,mid,i<<1);
    if(jr>mid) ans+=query(jl,jr,jt,mid+1,r,i<<1|1);
    return ans;
}
void build(int l,int r,int i){
    if(l==r){
        sum[i]=a[l];
        pf[i]=a[l]*a[l];
        
    }
    else {
    int mid=(l+r)>>1;
    build(l,mid,i<<1);
    build(mid+1,r,i<<1|1);
    up(i);
    }
    lazy[i]=0;
}
int main(){
    ios_base::sync_with_stdio(false);
    cin.tie(nullptr);
    int n,m;
    cin>>n>>m;
    for(int i=1;i<=n;i++){
        cin>>a[i];
    }
    build(1,n,1);
    while(m--){
        int opt;
        cin>>opt;
        if(opt==1){
            int l,r;
            double v;
            cin>>l>>r>>v;
            add(l,r,v,1,n,1);
        }else if(opt==2){
            int l,r;
            cin>>l>>r;
            printf("%.4lf\n",query(l,r,0,1,n,1)/(r-l+1));
        }else {
            int l,r;
            cin>>l>>r;
            printf("%.4lf\n",query(l,r,1,1,n,1)/(r-l+1)-query(l,r,0,1,n,1)/(r-l+1)*query(l,r,0,1,n,1)/(r-l+1));
        }
    }
}
\end{lstlisting}

\algo{Kruskal重构树}
\begin{lstlisting}
//Kruskal重构树: 最小生成树上任意两点路径最大值=新节点LCA的权值
const int maxn=1e5+10,maxm=2*maxn,maxk=3*maxn,maxh=20;
struct edge{ int u,v,w; }e[maxk];
int n,m,fa[maxk],head[maxk],tot=1;
struct graph{ int t,n; }g[maxk];
void add(int u,int v){ g[++tot].n=head[u]; g[tot].t=v; head[u]=tot; }
int nodekey[maxk],cntu,dep[maxk],st[maxk][maxh];
int find(int x){ return x==fa[x]?x:fa[x]=find(fa[x]); }
bool cmp(edge x,edge y){ return x.w<y.w; }
void kruskalbuild(){
    for(int i=1;i<=n;i++) fa[i]=i;
    sort(e+1,e+m+1,cmp); cntu=n;
    for(int i=1,fx,fy;i<=m;i++){
        fx=find(e[i].u); fy=find(e[i].v);
        if(fx!=fy){
            fa[fx]=fa[fy]=++cntu; fa[cntu]=cntu;
            nodekey[cntu]=e[i].w;
            add(cntu,fx); add(cntu,fy);
        }
    }
}
void dfs(int u,int f){
    st[u][0]=f; dep[u]=dep[f]+1;
    for(int p=1;p<maxh;p++) st[u][p]=st[st[u][p-1]][p-1];
    for(int e=head[u];e;e=g[e].n) dfs(g[e].t,u);
}
int lca(int u,int v){
    if(dep[u]<dep[v]) swap(u,v);
    for(int p=maxh-1;p>=0;p--) if(dep[st[u][p]]>=dep[v]) u=st[u][p];
    if(u==v) return u;
    for(int p=maxh-1;p>=0;p--) if(st[u][p]!=st[v][p]){ u=st[u][p]; v=st[v][p]; }
    return st[u][0];
}
\end{lstlisting}

\algo{左偏树}
\begin{lstlisting}
//可并堆结构，两个堆的合并时间复杂度为log（n）
//查询操作需要并查集的优化
//P3377,模板,删除一个堆中最小值
#include<bits/stdc++.h>
using namespace std;
const int maxn=100001;
int n,m;
int num[maxn],lef[maxn],righ[maxn],dis[maxn];
int father[maxn];
//father并不是代表该节点的父亲节点，而是代表并查集所对应的路径信息
//需要保证，并查集最上方代表节点=左偏树的头，也就是堆顶
void prepare(){
    dis[0]=-1;
    for(int i=1;i<=n;i++){
        lef[i]=0;
        righ[i]=0;
        dis[i]=0;
        father[i]=i;
    }
}
int find(int i){
    father[i]=father[i]==i?i:find(father[i]);
    return father[i];
}
int merge(int i,int j){
    if(i==0||j==0) return i+j;
    //维护小根堆，如果值一样，编号小的节点做头,大根堆同理
    if(num[i]>num[j]||(num[i]==num[j]&&i>j)){
        swap(i,j);
    }
    righ[i]=merge(righ[i],j);
    if(dis[lef[i]]<dis[righ[i]]){
        swap(lef[i],righ[i]);
    }
    dis[i]=dis[righ[i]]+1;
    father[lef[i]]=father[righ[i]]=i;
    return i;
}
//节点i一定是左偏树的头，在左偏树上删除节点i，返回新树的头结点编号
int pop(int i){
    father[lef[i]]=lef[i];
    father[righ[i]]=righ[i];
    father[i]=merge(lef[i],righ[i]);//保证删除节点后，能够正确再次查询
    lef[i]=righ[i]=dis[i]=0;
    return father[i];
}
int main(){
    cin>>n>>m;
    prepare();
    for(int i=1;i<=n;i++) cin>>num[i];
    for(int i=1;i<=m;i++){
        int op;
        cin>>op;
        if(op==1){
            int l,r;
            cin>>l>>r;
            if(num[l]!=-1&&num[r]!=-1){
                l=find(l),r=find(r);
                if(l!=r) merge(l,r);
            }
        }else {
            int x;
            cin>>x;
            if(num[x]==-1) cout<<-1<<endl;
            else {
            int ans=find(x);
            cout<<num[ans]<<endl;
            pop(ans);
            num[ans]=-1;
        }
    }
    }
}
\end{lstlisting}

\algo{虚树的构建}
\begin{lstlisting}
#include<bits/stdc++.h>
using namespace std;
const int maxn=1e5+10;
int h[maxn<<1],ne[maxn<<1],to[maxn<<1],tot1;//初始树
int n,q,k;
int hn[maxn<<1],nn[maxn<<1],tn[maxn<<1],tot2;//新树
int dep[maxn],st[maxn][20],dfn[maxn],dfncnt;
vector<int>arr,tmp;//关键点
bool iskey[maxn];
int cost[maxn],siz[maxn];
bool cmp(int x,int y){
    return dfn[x]<dfn[y];
}
void add1(int u,int v){
    ne[++tot1]=h[u];
    to[tot1]=v;
    h[u]=tot1;
}
void add2(int u,int v){
    nn[++tot2]=hn[u];
    tn[tot2]=v;
    hn[u]=tot2;
}
void dfs(int u,int f){
    dep[u]=dep[f]+1;
    dfn[u]=++dfncnt;
    st[u][0]=f;
    for(int i=1;i<20;i++){
        st[u][i]=st[st[u][i-1]][i-1];
    }
    for(int e=h[u];e;e=ne[e]){
        int v=to[e];
        if(v==f) continue;
        dfs(v,u);
    }
}
int getlca(int u,int v){
    if(dep[u]!=dep[v]){
        if(dep[u]<dep[v]) swap(u,v);
        for(int p=19;p>=0;p--){
            if(dep[st[u][p]]>=dep[v]){
                u=st[u][p];
            }
        }
    }
    if(u==v) return u;
    for(int p=19;p>=0;p--){
        if(st[u][p]!=st[v][p]){
            u=st[u][p];
            v=st[v][p];
        }
    }
    return st[u][0];
}
int buildvtree(){
    sort(arr.begin(),arr.end(),cmp);
    tmp=arr;
    for(int i=0;i<(int)arr.size()-1;i++){
        tmp.push_back(getlca(arr[i],arr[i+1]));
    }
    sort(tmp.begin(),tmp.end(),cmp);
    tmp.erase(unique(tmp.begin(),tmp.end()),tmp.end());
    tot2=0;
    for(int i=0;i<(int)tmp.size();i++){
        hn[tmp[i]]=0;
    }
    for(int i=0;i<tmp.size()-1;i++){
        add2(getlca(tmp[i],tmp[i+1]),tmp[i+1]);
    }
    return tmp[0];
}
\end{lstlisting}

\section{图论}

\algo{最短路}
\begin{lstlisting}
bitset<N> b[N];
for (int j = 1; j <= n; ++j) {  // 注意中转点在最外层
    for (int i = 1; i <= n; ++i) {
        if (b[i][j]) { b[i] |= b[j]; }
    }
}
\end{lstlisting}

\algo{最小生成树}
\begin{lstlisting}
int n, m;
vector<pii> G[N];
int kruskal()
{   // 默认是一张连通图
	vector <array<int,3>> e(m);
	for (auto& [w, u, v]: e) { cin >> u >> v >> w; }
	sort(e.begin(), e.end());
	int ans = 0, cnt = 0;
	d = dsu(n);  // 见并查集模板
	for (auto& [w, u, v]: e) {
		if (d.find(u) != d.find(v)) {
			d.merge(u, v);
			ans += w;
			++cnt;
			// 以下两行为建图，完成后 G 中存放了最小生成树
			G[u].push_back( { v, w } );
			G[v].push_back( { u, w } );
		}
		if (cnt == n-1) { break; }
	}
	return ans;
}
\end{lstlisting}

\algo{拓扑排序}
\begin{lstlisting}
queue<int> q;
for (int i = 1; i <= n; ++i) {
	if (deg[i] == 0) { q.push(i); }
}
vector<int> topo;
while (!q.empty()) {
	int u = q.front(); q.pop();
	topo.push_back(u);
	for (int v: G[u]) {
		if (--deg[v] == 0) { q.push(v); }
	}
}
\end{lstlisting}

\algo{LCA}
\begin{lstlisting}
const int N=5e5+5, I=20;
int n, m, s, dep[N], f[I+5][N], mn[I+5][N];
vector <pii> G[N];

void dfs(int u, int fa)
{   // 预处理父节点、深度、路径上的最小值
    f[0][u] = fa;
    dep[u] = dep[fa] + 1;
    for (int i = 1; i <= I; ++i) {
        f[i][u] = f[i-1][ f[i-1][u] ];
        mn[i][u] = min( mn[i-1][u], mn[i-1][ f[i-1][u] ] );
    }
    for (auto [v, w]: G[u]) {
        if (v != fa) {
            mn[0][v] = w;
            dfs(v, u);
        }
    }
}

int LCA(int u, int v)
{   // 求 u,v 到其 LCA 路径上的最小值
    int ans = numeric_limits<int>::max();
    function<void(int&,int)> jump = [&](int &u, int i) {
        ans = min(ans, mn[i][u]);
        u = f[i][u];
    };
    if (dep[u] < dep[v]) { swap(u, v); }
    for (int i = I; i >= 0; --i) {
        if (dep[f[i][u]] >= dep[v]) { jump(u, i); }
    }
    if (u == v) { return ans; }
    for (int i = I; i >= 0; --i) {
        if (f[i][u] != f[i][v]) {
            jump(u, i); jump(v, i);
        }
    }
    jump(u, 0); jump(v, 0);
    return ans;
}
\end{lstlisting}

\algo{树的直径与中心}
\begin{lstlisting}
int d[N], mxd;
void dfs(int u, int fa)
{
	for (auto v: G[u]) {
		if (v != fa) {
			dfs(v, u);
			mxd = max(mxd, d[u]+d[v]+1);
			d[u] = max(d[u], d[v]+1);
		}
	}
}
\end{lstlisting}

\algo{点分治与树的重心}
\begin{lstlisting}
int n, siz[N], dis[N];
bool vis[N];
vector<pii> G[N];

void calcsiz(int u, int fa, int sum, int &core)
{   // 求树的重心，sum 表述当前树的大小
    siz[u] = 1;
    int mxz = 0;
    for (auto [v, w]: G[u]) {
        if (v != fa && !vis[v]) {
            calcsiz(v, u, sum, core);
            siz[u] += siz[v];
            mxz = max(mxz, siz[v]);
        }
    }
    mxz = max(mxz, sum-siz[u]);
    if (mxz * 2 <= sum) { core = u; }
}

void calcdis(int u, int fa) 
{
    for (auto [v, w]: G[u]) {
        if (v != fa && !vis[v]) {
            dis[v] = dis[u] + w;
            calcdis(v, u);
        }
    }
}

void work(int u, int fa)
{
    vis[u] = true;
    for (auto [v, w]: G[u]) {
        if (v != fa && !vis[v]) {
            dis[v] = w;
            calcdis(v, u);
        }
    }
    for (auto [v, w]: G[u]) {
        if (v != fa && !vis[v]) {
            int sum = siz[v];
            int core = 0;
            calcsiz(v, u, sum, core);
            calcsiz(core, u, sum, core);  // 第二次调用更新 siz[]
            work(core, u);
        }
    }
}

int main()
{
    int core = 0;
    calcsiz(1, 0, n, core);
    calcsiz(core, 0, n, core);
    work(core, 0);
}
\end{lstlisting}

\algo{基环树}
\begin{lstlisting}
vector<int> G[N], cycle;
void dfs(int u, int fa)
{
    cycle.push_back(u);
    if (u == t) {
        sort(cycle.begin(), cycle.end());
        for (int x: cycle) { cout << x << " "; }
        exit(0);
    }
    for (int v: G[u]) {
        if (v != fa) { dfs(v, u); }
    }
    cycle.pop_back();
}
int main()
{
	dsu d(n);  // 见并查集板子
	for (int i = 1; i <= n; ++i) {
	    int u, v;
	    cin >> u >> v;
	    G[u].push_back(v);
	    G[v].push_back(u);
	    if (d.same(u, v)) { s=u; t=v; }
	    d.merge(u, v);
	}
	dfs(s, 0);
}
\end{lstlisting}

\algo{Hierholzer求欧拉路}
\begin{lstlisting}
int n, deg[N<<1], cur[N<<1];
bool vis[N];
vector<int> ans;
vector<pair<int,int>> G[N<<1];

void dfs(int u)
{
    for (int &j = cur[u]; j < G[u].size(); ++j) {
        auto [v, i] = G[u][j];
        if (!vis[i]) {
            vis[i] = true;
            dfs(v);
            ans.push_back(i);
        }
    }
}

int main()
{
	int s = 1;
	while (s <= 2*n && deg[s] % 2 == 0) { ++s; }
	if (s > 2 * n) {
	    s = 1;
	    while (!deg[s]) { ++s; }
	}
	
	int cnt = 0;
	for (int i = 1; i <= 2*n; ++i) {
	    cnt += deg[i] % 2;
	}
	if (cnt > 2) {
	    cout << "NO\n"; return;
	}
	
	dfs(s);
	
	if (ans.size() != n) {
	    cout << "NO\n";
	} else {
	    cout << "YES\n";
	    for (auto x: ans) {
	        cout << x << " \n"[x == ans.back()];
	    }
	}
}
\end{lstlisting}

\algo{DSU on tree}
\begin{lstlisting}
const int N=2e5+5;
int n, c[N], siz[N], hvs[N];
int L[N], R[N], dfn[N], totdfn;
vector<int> G[N];

void dfs_init(int u, int fa)
{   // 预处理 dfn 和重儿子
	siz[u] = 1;
	L[u] = ++totdfn;
	dfn[totdfn] = u;
	for (int v: G[u]) {
		if (v != fa) {
			dfs_init(v, u);
			siz[u] += siz[v];
			if (hvs[u] == 0 || siz[v] > siz[hvs[u]]) {
				hvs[u] = v;
			}
		}
	}
	R[u] = totdfn;
}

void add(int u, int dt)
{
	--ccnt[cnt[c[u]]];
	cnt[c[u]] += dt;
	++ccnt[cnt[c[u]]];
}

void dfs_solve(int u, int fa, bool keep)
{
	for (auto v: G[u]) {  // 先做轻儿子，不保留
		if (v != fa && v != hvs[u]) {
			dfs_solve(v, u, false);
		}
	}
	if (hvs[u]) {  // 再做重儿子，保留
		dfs_solve(hvs[u], u, true);
	}
	for (auto v: G[u]) {  // 再把轻儿子的信息加上
		if (v != fa && v != hvs[u]) {
			for (int i = L[v]; i <= R[v]; ++i) {
				add(dfn[i], 1);
			}
		}
	}
	add(u, 1);
	ans += (cnt[c[u]] * ccnt[cnt[c[u]]] == siz[u]);
	if (keep) { return; }
	for (int i = L[u]; i <= R[u]; ++i) {
		add(dfn[i], -1);
	}
}
\end{lstlisting}

\algo{tarjan}
\begin{lstlisting}
struct SCC {
	int n, totdfn=0, colcnt=0;
	stack<int> stk;
	vector<int> dfn, low, col;
	vector<vector<int>> G;
	SCC(int n): n(n) {
		G.assign(n+1, {});
		dfn.assign(n+1, 0);
		low.assign(n+1, 0);
		col.assign(n+1, 0);
	}
	void add_edge(int u, int v) {
		G[u].push_back(v);
	}
	void dfs(int u) {
		stk.push(u);
		dfn[u] = low[u] = ++totdfn;
		for (auto v: G[u]) {
			if (!dfn[v]) {
				dfs(v);
				low[u] = min(low[u], low[v]);
			} else if (!col[v]) {
				low[u] = min(low[u], dfn[v]);
			}
		}
		if (low[u] == dfn[u]) {
			col[u] = ++colcnt;
			int v;
			do {
				v = stk.top(); stk.pop();
				col[v] = colcnt;
			} while(v != u);
		}
	}
	vector<int> work() {
		for (int i = 1; i <= n; ++i) {
			if (!dfn[i]) { dfs(i); }
		}
		return col;
	}
};
\end{lstlisting}

\algo{树哈希}
\begin{lstlisting}
using ull = unsigned long long;
const ull msk = chrono::steady_clock::now().time_since_epoch().count();

// ull h(ull x) { return x * x * x * 1237123 + 19260817; }
// ull f(ull x) { return h(x & ((1 << 31) - 1)) + h(x >> 31); }
ull shift(ull x)
{
    x ^= msk;
    x ^= x << 13;
    x ^= x >> 7;
    x ^= x << 17;
    x ^= msk;
    return x;
}

ull has[N];
void getHash(int u, int fa) 
{
    has[u] = 1;
    for (auto v: G[u]) {
        if (v == fa) { continue; }
        getHash(v, u);
        has[u] += shift(has[v]);
    }
}
\end{lstlisting}

\algo{Bellman-Ford / SPFA（判负环）}
\begin{lstlisting}
//Bellman-Ford + SPFA优化, O(n*m), 可用于判负环
#include<bits/stdc++.h>
using namespace std; 
const int maxn=2010,maxm=3010,inf=1e8;
int head[maxm<<1],n,m,cnt,upfill[maxn],dis[maxn];
queue<int>q; bool enter[maxn];
struct graph{ int next,to,w; }g[maxm<<1];
void add(int u,int v,int w){ g[cnt].next=head[u]; g[cnt].to=v; g[cnt].w=w; head[u]=cnt++; }
void build(int n){ cnt=1; while(!q.empty())q.pop();
    for(int i=1;i<=n;i++){ head[i]=0; upfill[i]=0; dis[i]=inf; enter[i]=0; } }
bool spfa(){ dis[1]=0; upfill[1]++; q.push(1); enter[1]=true;
    while(!q.empty()){ int x=q.front(); q.pop(); enter[x]=false;
        for(int ei=head[x];ei;ei=g[ei].next){ int v=g[ei].to,w=g[ei].w;
            if(dis[x]+w<dis[v]){ dis[v]=dis[x]+w;
                if(!enter[v]){ if(upfill[v]++==n)return true; enter[v]=true; q.push(v); } } } }
    return false; }
\end{lstlisting}

\algo{Floyd 全源最短路}
\begin{lstlisting}
//全源最短路 O(n^3), 允许负边但不能有负环
const int maxn=1001,inf=1e8; int dis[maxn][maxn];
void floyd(){ for(int b=1;b<=n;b++)for(int i=1;i<=n;i++)for(int j=1;j<=n;j++)
    if(dis[i][b]!=inf&&dis[b][j]!=inf&&dis[i][j]>dis[i][b]+dis[b][j])
        dis[i][j]=dis[i][b]+dis[b][j]; }
\end{lstlisting}

\algo{树上倍增}
\begin{lstlisting}
//快速查询从节点往上走k步的祖先, 预处理O(n log n), 查询O(log n)
const int maxn=5e5+10; int st[maxn][32],deep[maxn],power;
void dfs(int i,int f){ deep[i]=(i==0?1:deep[f]+1); st[i][0]=f;
    for(int p=1;(1<<p)<=deep[i];p++) st[i][p]=st[st[i][p-1]][p-1];
    for(int ei=head[i];ei;ei=tr[ei].n) dfs(tr[ei].t,i); }
int getkthAncestor(int i,int k){ if(deep[i]<=k)return -1; int s=deep[i]-k;
    for(int p=power;p>=0;p--) if(deep[st[i][p]]>=s) i=st[i][p]; return i; }
\end{lstlisting}

\algo{树的直径}
\begin{lstlisting}
//两次DFS求树的直径
int maxd=0,far=1;
void dfs(int u,int fa,int d){ if(d>maxd){maxd=d;far=u;} for(int v:g[u])if(v!=fa)dfs(v,u,d+1); }
// dfs(1,0,0); dfs(far,0,0); maxd即为直径
\end{lstlisting}

\algo{最大流}
\begin{lstlisting}
const ll INF = 9e18; 
struct Flow {
    vector<tuple<int,ll,int>> G[N];  // 终点，流量，反边
    int dis[N], now[N];  // now[u] 为当前弧
    void add_edge(int u, int v, ll w) {
        G[u].push_back( { v, w, G[v].size() } );
        G[v].push_back( { u, 0, G[u].size()-1 } );
    }
    void bfs(int st) {  // 求层次图
        queue<int> q;
        memset(dis, 0, sizeof(dis));
        q.push(st); dis[st] = 1;
        while (!q.empty()) {
            int u = q.front(); q.pop();
            for (auto& [v, w, inv]: G[u]) {
                if (w != 0 && dis[v] == 0) {
                    dis[v] = dis[u] + 1;
                    q.push(v);
                }
            }
        }
    }
    ll dfs(int u, int t, ll flow) {  // 求阻塞流
        if (u == t) { return flow; }
        for (int &i=now[u]; i < (int)G[u].size(); ++i) {
            auto& [v, w, inv] = G[u][i];
            if (w != 0 && dis[v] > dis[u]) {
                ll d = dfs(v, t, min(flow, w));
                if (d > 0) {
                    w -= d;
                    get<1>(G[v][inv]) += d;
                    return d;
                }
            }
        }
        return 0;
    }
    ll dicnic(int st, int ed) {
        for (ll flow=0, res;;) {
            bfs(st);
            if (dis[ed] == 0) { return flow; }
            memset(now, 0, sizeof(now));
            while ((res=dfs(st,ed,INF)) > 0) {
                flow += res;
            }
        }
    }
};
\end{lstlisting}

\algo{最近公共祖先（LCA）}
\begin{lstlisting}
//利用树链剖分的方法
#include<bits/stdc++.h>
using namespace std;
const int N=5e5+5;
struct node{
	int next,to;
}e[N<<1];
int tot,head[N];
int dep[N],siz[N],hson[N],fa[N];
int top[N];
int n,m,s;
void add(int x,int y){
	e[++tot].to=y;
	e[tot].next=head[x];
	head[x]=tot;
}
void dfs1(int x,int f){
	siz[x]=1;
	fa[x]=f;
	dep[x]=dep[f]+1;
	for(int i=head[x];i;i=e[i].next){
		int y=e[i].to;
		if(y!=f){
			dfs1(y,x);
			siz[x]+=siz[y];
			if(siz[hson[x]]<siz[y] || !hson[x]){
				hson[x]=y;
			}
		}
	}
}
void dfs2(int x,int t){
	top[x]=t;
	if(!hson[x])return ;
	dfs2(hson[x],t);
	for(int i=head[x];i;i=e[i].next){
		int y=e[i].to;
		if(y!=fa[x] && y!=hson[x]){
			dfs2(y,y);
		}
	}
}
int LCA(int x,int y){
	while(top[x]!=top[y]){
		if(dep[top[x]]<dep[top[y]]){
			swap(x,y);
		}
		x=fa[top[x]];
	}
	return dep[x]<dep[y]?x:y;
}
int main(){
	ios::sync_with_stdio(false);
	cin.tie(NULL);
	cout.tie(NULL);
	cin>>n>>m>>s;
	for(int i=1;i<n;i++){
		int x,y;
		cin>>x>>y;
		add(x,y);
		add(y,x);
	}
	dfs1(s,0);
	dfs2(s,s);
	for(int i=1;i<=m;i++){
		int a,b;
		cin>>a>>b;
		cout<<LCA(a,b)<<'\n';
	}
	return 0;
}
\end{lstlisting}

\algo{树链剖分}
\begin{lstlisting}
#include<algorithm>
#include<iostream>
#include<cstdlib>
#include<cstring>
#include<cstdio>
#define Rint register int
#define mem(a,b) memset(a,(b),sizeof(a))
#define Temp template<typename T>
using namespace std;
typedef long long LL;
Temp inline void read(T &x){
    x=0;T w=1,ch=getchar();
    while(!isdigit(ch)&&ch!='-')ch=getchar();
    if(ch=='-')w=-1,ch=getchar();
    while(isdigit(ch))x=(x<<3)+(x<<1)+(ch^'0'),ch=getchar();
    x=x*w;
}

#define mid ((l+r)>>1)
#define lson rt<<1,l,mid
#define rson rt<<1|1,mid+1,r
#define len (r-l+1)

const int maxn=200000+10;
int n,m,r,mod;
//见题意 
int e,beg[maxn],nex[maxn],to[maxn],w[maxn],wt[maxn];
//链式前向星数组，w[]、wt[]初始点权数组 
int a[maxn<<2],laz[maxn<<2];
//线段树数组、lazy操作 
int son[maxn],id[maxn],fa[maxn],cnt,dep[maxn],siz[maxn],top[maxn]; 
//son[]重儿子编号,id[]新编号,fa[]父亲节点,cnt dfs_clock/dfs序,dep[]深度,siz[]子树大小,top[]当前链顶端节点 
int res=0;
//查询答案 

inline void add(int x,int y){//链式前向星加边 
    to[++e]=y;
    nex[e]=beg[x];
    beg[x]=e;
}
//-------------------------------------- 以下为线段树 
inline void pushdown(int rt,int lenn){
	laz[rt<<1]+=laz[rt];
	laz[rt<<1|1]+=laz[rt];
    a[rt<<1]+=laz[rt]*(lenn-(lenn>>1));
    a[rt<<1|1]+=laz[rt]*(lenn>>1);
    a[rt<<1]%=mod;
    a[rt<<1|1]%=mod;
    laz[rt]=0;
}

inline void build(int rt,int l,int r){
    if(l==r){
        a[rt]=wt[l];
        if(a[rt]>mod)a[rt]%=mod;
        return;
    }
    build(lson);
    build(rson);
    a[rt]=(a[rt<<1]+a[rt<<1|1])%mod;
}

inline void query(int rt,int l,int r,int L,int R){
	if(L<=l&&r<=R){res+=a[rt];res%=mod;return;}
	else{
		if(laz[rt])pushdown(rt,len);
		if(L<=mid)query(lson,L,R);
		if(R>mid)query(rson,L,R);
	}
}

inline void update(int rt,int l,int r,int L,int R,int k){
    if(L<=l&&r<=R){
		laz[rt]+=k;
		a[rt]+=k*len;
	}
	else{
		if(laz[rt])pushdown(rt,len);
		if(L<=mid)update(lson,L,R,k);
		if(R>mid)update(rson,L,R,k);
		a[rt]=(a[rt<<1]+a[rt<<1|1])%mod;
	}
}
//---------------------------------以上为线段树 
inline int qRange(int x,int y){
	int ans=0;
	while(top[x]!=top[y]){//当两个点不在同一条链上 
		if(dep[top[x]]<dep[top[y]])swap(x,y);//把x点改为所在链顶端的深度更深的那个点
		res=0;
		query(1,1,n,id[top[x]],id[x]);//ans加上x点到x所在链顶端 这一段区间的点权和
		ans+=res;
		ans%=mod;//按题意取模 
		x=fa[top[x]];//把x跳到x所在链顶端的那个点的上面一个点
	}
	//直到两个点处于一条链上
	if(dep[x]>dep[y])swap(x,y);//把x点深度更深的那个点
	res=0;
	query(1,1,n,id[x],id[y]);//这时再加上此时两个点的区间和即可
	ans+=res;
	return ans%mod;
}

inline void updRange(int x,int y,int k){//同上 
	k%=mod;
	while(top[x]!=top[y]){
		if(dep[top[x]]<dep[top[y]])swap(x,y);
		update(1,1,n,id[top[x]],id[x],k);
		x=fa[top[x]];
	}
	if(dep[x]>dep[y])swap(x,y);
	update(1,1,n,id[x],id[y],k);
}

inline int qSon(int x){
	res=0;
	query(1,1,n,id[x],id[x]+siz[x]-1);//子树区间右端点为id[x]+siz[x]-1 
	return res;
}

inline void updSon(int x,int k){//同上 
	update(1,1,n,id[x],id[x]+siz[x]-1,k);
}

inline void dfs1(int x,int f,int deep){//x当前节点，f父亲，deep深度 
	dep[x]=deep;//标记每个点的深度 
	fa[x]=f;//标记每个点的父亲 
	siz[x]=1;//标记每个非叶子节点的子树大小 
	int maxson=-1;//记录重儿子的儿子数 
	for(Rint i=beg[x];i;i=nex[i]){
		int y=to[i];
		if(y==f)continue;//若为父亲则continue 
		dfs1(y,x,deep+1);//dfs其儿子 
		siz[x]+=siz[y];//把它的儿子数加到它身上 
		if(siz[y]>maxson)son[x]=y,maxson=siz[y];//标记每个非叶子节点的重儿子编号 
	}
}

inline void dfs2(int x,int topf){//x当前节点，topf当前链的最顶端的节点 
	id[x]=++cnt;//标记每个点的新编号 
	wt[cnt]=w[x];//把每个点的初始值赋到新编号上来 
	top[x]=topf;//这个点所在链的顶端 
	if(!son[x])return;//如果没有儿子则返回 
	dfs2(son[x],topf);//按先处理重儿子，再处理轻儿子的顺序递归处理 
	for(Rint i=beg[x];i;i=nex[i]){
		int y=to[i];
		if(y==fa[x]||y==son[x])continue;
		dfs2(y,y);//对于每一个轻儿子都有一条从它自己开始的链 
	}
}

int main(){
    read(n);read(m);read(r);read(mod);
    for(Rint i=1;i<=n;i++)read(w[i]);
    for(Rint i=1;i<n;i++){
        int a,b;
        read(a);read(b);
        add(a,b);add(b,a);
    }
    dfs1(r,0,1);
    dfs2(r,r);
    build(1,1,n);
	while(m--){
        int k,x,y,z;
        read(k);
        if(k==1){
            read(x);read(y);read(z);
            updRange(x,y,z);
        }
        else if(k==2){
            read(x);read(y);
            printf("%d\n",qRange(x,y));
        }
        else if(k==3){
            read(x);read(y);
            updSon(x,y);
        }
        else{
            read(x);
            printf("%d\n",qSon(x));
    	}
    }
}
\end{lstlisting}

\algo{最大路匹配}
\begin{lstlisting}
#include<bits/stdc++.h>
using namespace std;
const int N=510;
vector<int>g[N];
bool vis[N];
int match[N];
bool dfs(int u){
    for(int v:g[u]){
        if(!vis[v]){
            vis[v]=1;
            if(match[v]==0||dfs(match[v])){
                match[v]=u;
                return true;
            }
        }
    }
    return false;
}
int main(){
    int n,m,ans=0;
    cin>>n>>m;
    for(int i=1;i<=n;i++){
        int x;
        cin>>x;
        for(int j=1;j<=x;j++){
            int y;
            cin>>y;
            g[i].push_back(y);
        }
    }
    for(int i=1;i<=n;i++){
        memset(vis,0,sizeof(vis));
        if(dfs(i)) ans++;
    }
    cout<<ans;
}
\end{lstlisting}

\algo{Dijkstra算法}
\begin{lstlisting}
//小根堆实现较为简单，这里采用以点为复杂度的反向索引堆的优化,一般不会用
#include<bits/stdc++.h>
using namespace std;
const int maxn=100010;
const int inf=1e8;
struct graph{
    int next,to,w;
}g[maxn*2];
int head[maxn],cnt,n,m;
int heap[maxn];//反向索引堆
int where[maxn];//表示点在堆上位置，初始化为-1
int heapSize;
int dis[maxn];
void build(int n){
    cnt=1;
    heapSize=0;
    for(int i=1;i<=n;i++){
        heap[i]=0;
        where[i]=-1;
        head[i]=0;
        dis[i]=inf;
    }
}
void add(int u,int v,int w){
    g[cnt].next=head[u];
    g[cnt].to=v;
    g[cnt].w=w;
    head[u]=cnt++;
}
void Swap(int i,int j){
    swap(heap[i],heap[j]);
    where[heap[i]]=i;
    where[heap[j]]=j;
}
void heapInsert(int i){
    while(dis[heap[i]]<dis[heap[(i-1)/2]]){
        Swap(i,(i-1)/2);
        i=(i-1)/2;
    }
}
void addorupdateoringnore(int v,int c){
    if(where[v]==-1){
        heap[heapSize]=v;
        where[v]=heapSize++;
        dis[v]=c;
        heapInsert(where[v]);
    }else if(where[v]>=0){
        dis[v]=min(dis[v],c);
        heapInsert(where[v]);
    }
}
void heapify(int i){
    int l=1;
    while(l<heapSize){
        int best=l+1<heapSize&&dis[heap[l+1]]<dis[heap[l]]?l+1:l;
        best=dis[heap[best]]<dis[heap[i]]?best:i;
        if(best==i) break;
        Swap(best,i);
        i=best;
        l=i*2+1;
    }
}
bool isEmpty(){
    return heapSize==0;
}
int pop(){
    int ans=heap[0];
    swap(heap[0],heap[--heapSize]);
    heapify(0);
    where[ans]=-2;
    return ans;
}
int main(){
    cin>>n>>m;
    int s;
    cin>>s;
    build(n);
    for(int i=1;i<=m;i++){
        int u,v,w;
        cin>>u>>v>>w;
        add(u,v,w);
    }
    addorupdateoringnore(s,0);
    while(!isEmpty()){
        int u=pop();
        for(int ei=head[u];ei>0;ei=g[ei].next){
            addorupdateoringnore(g[ei].to,dis[u]+g[ei].w);
        }
    }
    int ans=-inf;
    for(int i=1;i<=n;i++){
        /*if(dis[i]==inf){
            cout<<-1<<endl;
        }
        ans=max(ans,dis[i]);*/
        cout<<dis[i]<<' ';
    }
    //cout<<ans<<endl;
}
\end{lstlisting}

\algo{分层图最短路}
\begin{lstlisting}
//无论是在bfs上分层还是dijkstra上，本质其实就是把原图上的点与其的不同状态整合出更多的点
//难点在于如何将点上的状态给储存下来，以及如何转移
#include<bits/stdc++.h>
using namespace std;
const int maxn=50010;
const int inf=1e8;
struct dot{
    int x,cost,k;
};
struct graph{
    int next,to,w;
}g[maxn<<1];
struct compare{
   bool operator()(dot x,dot y){
        return x.cost>y.cost;
        // 小根堆
    }
};
int head[maxn<<1],cnt=1,n,m,k,dis[maxn][20];
bool vis[maxn][20];
priority_queue<dot,vector<dot>,compare>q;
void add(int u,int v,int w){
    g[cnt].next=head[u];
    g[cnt].to=v;
    g[cnt].w=w;
    head[u]=cnt++;
}
int main(){
    cin>>n>>m>>k;
    int s,t;
    cin>>s>>t;
    for(int i=1;i<=m;i++){
        int u,v,w;
        cin>>u>>v>>w;
        add(u,v,w);
        add(v,u,w);
    }
    for(int i=0;i<=n;i++){
        for(int j=0;j<=k;j++){
            dis[i][j]=inf;
        }
    }
    dis[s][0]=0;
    q.push({s,0,0});
    while(!q.empty()){
        int xn=q.top().x;
        int cost=q.top().cost;
        int kn=q.top().k;
        q.pop();
        if(vis[xn][kn]) continue;
        vis[xn][kn]=1;
        if(xn==t){
            cout<<cost<<endl;
            return 0;
        }
        for(int i=head[xn];i>0;i=g[i].next){
            int v=g[i].to;
            if(kn<k&&cost<dis[v][kn+1]){
                dis[v][kn+1]=cost;
                q.push({v,cost,kn+1});
            }
            if(cost+g[i].w<dis[v][kn]){
                dis[v][kn]=cost+g[i].w;
                q.push({v,cost+g[i].w,kn});
            }
        }
    }   
}
/*
struct Compare {
    bool operator()(const State& a, const State& b) {
        return a.cost > b.cost; // 小根堆
    }
};
    priority_queue<State, vector<State>, Compare> q;
    这样子定义也可以
*/
\end{lstlisting}

\algo{双向广搜}
\begin{lstlisting}
#include<bits/stdc++.h>
using namespace std;
typedef long long ll;
ll n,m;
ll an[45];
ll l[1<<20+1],r[1<<20+1];
ll f(int i,int e,ll s,ll w,ll a[],ll j){
    if(s>w) return j;
    if(i==e) a[j++]=s;
    else {
        j=f(i+1,e,s,w,a,j);
        j=f(i+1,e,s+an[i],w,a,j);
    }
    return j;
}
ll compute(){
    ll lsize=f(0,n>>1,0,m,l,0);
    ll rsize=f(n>>1,n,0,m,r,0);
    sort(l,l+lsize);
    sort(r,r+rsize);
    ll ans=0;
    for(int i=lsize-1,j=0;i>=0;i--){
        while(j<rsize&&l[i]+r[j]<=m){
            j++;
        }
        ans+=j;
    }
    return ans;
}
int main(){
    cin>>n>>m;
  for(int i=0;i<n;i++) cin>>an[i];
  sort(an,an+n);
  cout<<compute();
}
\end{lstlisting}

\algo{tarjan(强连通分量)}
\begin{lstlisting}
//强连通分量求解
int dfn[N], low[N], dfncnt, s[N], in_stack[N], tp;
int scc[N], sc;  // 结点 i 所在 SCC 的编号
int sz[N];       // 强连通 i 的大小

void tarjan(int u) {
  low[u] = dfn[u] = ++dfncnt, s[++tp] = u, in_stack[u] = 1;
  for (int i = h[u]; i; i = e[i].nex) {
    const int &v = e[i].t;
    if (!dfn[v]) {
      tarjan(v);
      low[u] = min(low[u], low[v]);
    } else if (in_stack[v]) {
      low[u] = min(low[u], dfn[v]);
    }
  }
  if (dfn[u] == low[u]) {
    ++sc;
    do {
      scc[s[tp]] = sc;
      sz[sc]++;
      in_stack[s[tp]] = 0;
    } while (s[tp--] != u);
  }
}
\end{lstlisting}

\algo{边双连通分量}
\begin{lstlisting}
#include <algorithm>
#include <iostream>
#include <stack>
#include <vector>

using namespace std;
constexpr int N = 5e5 + 5, M = 2e6 + 5;
int n, m;

struct edge {
  int to, nt;
} e[M << 2];

int hd[N << 1], tot = 1;

void add(int u, int v) { e[++tot] = edge{v, hd[u]}, hd[u] = tot; }

void uadd(int u, int v) { add(u, v), add(v, u); }

int bfncnt, sum;
int dfn[N], low[N],bcc[N],bcc_cnt;
bool vis[N];
vector<vector<int>> ans;
stack<int> st;
void tarjan(int u, int in) {
  low[u] = dfn[u] = ++dfncnt;
  st.push(u), vis[u] = true;
  for (int i = hd[u]; i; i = e[i].nt) {
    int v = e[i].to;
    if (i == (in ^ 1)) continue;
    if (!dfn[v])
      tarjan(v, i), low[u] = min(low[u], low[v]);
    else if (vis[v])
      low[u] = min(low[u], dfn[v]);
  }
  if (dfn[u] == low[u]) {
    bcc_cnt++;
    bcc[u]=bcc_cnt;
    vector<int> t;
    t.push_back(u);
    while (st.top() != u)
      t.push_back(st.top()), bcc[st.top()]=bcc_cnt,vis[st.top()] = false, st.pop();
    st.pop(), ans.push_back(t);
  }
}
\end{lstlisting}

\algo{割边}
\begin{lstlisting}
void tarjan(int u,int pre) {
  low[u] = dfn[u] = ++dfncnt;
  for (int ei = head[u]; ei; ei = e[ei].nex) {
    if((ei^1)==pre) continue;
    const int &v = e[ei].t;
    if (!dfn[v]) {
      tarjan(v,ei);
      
      low[u] = min(low[u], low[v]);
      if(low[v]>dfn[u]){
        cutedge[ei>>1]=1;//无向图割边
        ans1++;
      }
    } else  {
      low[u] = min(low[u], dfn[v]);
    }
  }
}
\end{lstlisting}

\algo{点双连通分量}
\begin{lstlisting}
void tarjan(int u) {
  dfn[u] = low[u] = ++bcc_cnt, sta[++top] = u;
  if (u == root && hd[u] == 0) { //root=i;
    dcc[++cnt].push_back(u);
    return;
  }
  int f = 0;
  for (int i = hd[u]; i; i = e[i].nt) {
    int v = e[i].to;
    if (!dfn[v]) {
      tarjan(v);
      low[u] = min(low[u], low[v]);
      if (low[v] >= dfn[u]) {
        if (++f > 1 || u != root) cut[u] = true;
        cnt++;
        do dcc[cnt].push_back(sta[top--]);
        while (sta[top + 1] != v);
        dcc[cnt].push_back(u);
      }
    } else
      low[u] = min(low[u], dfn[v]);
  }
}
\end{lstlisting}

\algo{割点}
\begin{lstlisting}
/*
洛谷 P3388 【模板】割点（割顶）
*/
#include <iostream>
#include <vector>
using namespace std;
int n, m;  // n：点数 m：边数
int dfn[100001], low[100001], idx, res;
// dfn：记录每个点的时间戳
// low：能不经过父亲到达最小的编号，idx：时间戳，res：答案数量
bool vis[100001], flag[100001];  // flag: 答案 vis：标记是否重复
vector<int> edge[100001];        // 存图用的

void Tarjan(int u, int fa) {  // u 当前点的编号，fa 自己爸爸的编号
  vis[u] = true;              // 标记
  low[u] = dfn[u] = ++idx;    // 打上时间戳
  int child = 0;              // 每一个点儿子数量
  for (const auto &v : edge[u]) {  // 访问这个点的所有邻居 （C++11）
    if (!vis[v]) {
      child++;                       // 多了一个儿子
      Tarjan(v, u);                  // 继续
      low[u] = min(low[u], low[v]);  // 更新能到的最小节点编号
      if (fa != u && low[v] >= dfn[u] && !flag[u]) {  // 主要代码
        // 如果不是自己，且不通过父亲返回的最小点符合割点的要求，并且没有被标记过
        // 要求即为：删了父亲连不上去了，即为最多连到父亲
        flag[u] = true;
        res++;  // 记录答案
      }
    } else if (v != fa) {
      // 如果这个点不是自己的父亲，更新能到的最小节点编号
      low[u] = min(low[u], dfn[v]);
    }
  }
  // 主要代码，自己的话需要 2 个儿子才可以
  if (fa == u && child >= 2 && !flag[u]) {
    flag[u] = true;
    res++;  // 记录答案
  }
}

int main() {
  cin >> n >> m;                  // 读入数据
  for (int i = 1; i <= m; i++) {  // 注意点是从 1 开始的
    int x, y;
    cin >> x >> y;
    edge[x].push_back(y);
    edge[y].push_back(x);
  }  // 使用 vector 存图
  for (int i = 1; i <= n; i++)  // 因为 Tarjan 图不一定连通
    if (!vis[i]) {
      idx = 0;       // 时间戳初始为 0
      Tarjan(i, i);  // 从第 i 个点开始，父亲为自己
    }
  cout << res << endl;
  for (int i = 1; i <= n; i++)
    if (flag[i]) cout << i << " ";  // 输出结果
  return 0;
}
\end{lstlisting}

\section{数学}

\algo{线性筛}
\begin{lstlisting}
vector<int> sieve(int n)
{
	vector<bool> vis(n+1);
    vector<int> mP(n+1,1);
	vector<int> prime;
	vis[1] = 1;
	for (int i = 2; i <= n; ++i) {
		if (!vis[i]) {
            prime.push_back(i); 
            mP[i] = i;
        }
		for (auto p: prime) {
			if (i*p > n) { break; }
			vis[i*p] = 1;
            mP[i*p] = p;
			if (i%p == 0) { break; }  // 线性筛关键优化
		}
	}
	return prime;
}

// 分解质因数
map<int,int> cnt;
for (int y = x; y > 1; y /= mP[y]) {
    ++cnt[mP[y]];
}
\end{lstlisting}

\algo{因数}
\begin{lstlisting}
const int X=1e5;
vector<int> fac[X+5];
for (int i = 1; i <= X; ++i) {
    for (int j = i; j <= X; j += i) {
        fac[j].push_back(i);
    }
}
\end{lstlisting}

\algo{快速幂与乘法逆元}
\begin{lstlisting}
const int MOD=998244353;
ll qpow(ll a, ll b=MOD-2)
{   // 快速幂函数，x 的乘法逆元为 qpow(x)
    ll res=1;
    for (; b; b>>=1) {
        if (b & 1) { res=(res*a)%MOD; }
        a = a*a%MOD;  // 不要少 MOD
    }
    return res;
}
\end{lstlisting}

\algo{排列组合数}
\begin{lstlisting}
const int N=1e5;
ll fac[N+5], Inv[N+5];
ll C(ll n, ll k) { return n==0? 1: fac[n] * Inv[k] % MOD * Inv[n-k] % MOD; }
ll P(ll n, ll k) { return n==0? 1: fac[n] * Inv[n-k] % MOD; }
// 预处理阶乘和阶乘的逆元
for (int i = 1; i <= N; ++i) {
	fac[i] = fac[i-1] * i % MOD;
	Inv[i] = qpow(fac[i]); 
}
\end{lstlisting}

\algo{Lucas}
\begin{lstlisting}
const int MOD = 1000003;
ll C(ll a, ll b) {
	if (a < b) { return 0; }
	ll down = 1, up = 1;
	for (int i = a, j = 1; j <= b; --i, ++j) {
		up = up * i % MOD;
		down = down * j % MOD;
	}
	return up * qpow(down) % MOD;
}
ll lucas(ll a, ll b) {
	if (a < p && b < p) { return C(a, b); }
	return C(a%MOD, b%MOD) * lucas(a/MOD, b/MOD) % MOD;
}
\end{lstlisting}

\algo{卡特兰数}
\begin{lstlisting}
ll Cat(ll n) { return C(2*n, n) * qpow(n+1) % MOD; }
\end{lstlisting}

\algo{数论分块}
\begin{lstlisting}
ll l=1, r=0;
while (l <= n) {
    r = n / (n/l);
    // 统计 ans
    l = r + 1;
}
\end{lstlisting}

\algo{高斯消元}
\begin{lstlisting}
int Guess(vector<vector<db>> &a)
{
	int n = a.size();
    int c, r;
    for (c = 0, r = 0; c < n; ++c) {
        // 选主元
        int t = r;
        for (int i = r + 1; i < n; ++i) {
            if (fabs(a[i][c]) > fabs(a[t][c])) { t = i; }
        }
        if (sgn(a[t][c]) == 0) { continue; }

        // 行交换
        for (int i = c; i <= n; ++i) { swap(a[t][i], a[r][i]); }

        // 归一化
        for (int i = n; i >= c; --i) { a[r][i] /= a[r][c]; }

        // 前向消元
        for (int i = r + 1; i < n; ++i) {
            if (sgn(a[i][c]) != 0) {
                for (int j = n; j >= c; --j) {
                    a[i][j] -= a[r][j] * a[i][c];
                }
            }
        }
        ++r;
    }

    if (r < n) {
        for (int i = r; i < n; ++i) {
            if (sgn(a[i][n]) != 0) { return 2; }  // 无解
        }
        return 1;  // 无穷多解
    }

    // 回代
    for (int i = n - 1; i >= 0; --i) {
        for (int j = i + 1; j < n; ++j) {
            a[i][n] -= a[i][j] * a[j][n];
        }
    }
    return 0;
}
\end{lstlisting}

\algo{矩阵快速幂}
\begin{lstlisting}
using matrix = vector<vector<ll>>;
matrix IM(int n) {
    matrix I(n, vector<ll>(n));
    for (int i = 0; i < n; ++i) { I[i][i] = 1; }
    return I;
}
matrix operator* (matrix A, matrix B) {
    assert(A.front().size() == B.size());
    int n = A.size(), m = B.size(), k = B.front().size();
    matrix C(n, vector<ll>(m));
    for (int i = 0; i < n; ++i) {
        for (int j = 0; j < k; ++j) {
            for (int x = 0; x < m; ++x) {
                C[i][j] = add(C[i][j], mul(A[i][x], B[x][j]));
            }
        }
    }
    return C;
}
matrix qpow(matrix a, ll b) {
    matrix res = IM(a.size());
    for (; b; b>>=1, a=a*a) {
        if (b & 1) { res = res * a; }
    }
    return res;
}
\end{lstlisting}

\algo{自适应辛普森法}
\begin{lstlisting}
double simpson(double l, double r) {
    double mid = (l + r) / 2;
    return (r - l) * (f(l) + 4 * f(mid) + f(r)) / 6;  // 辛普森公式
}

double asr(double l, double r, double eps, double ans,
           int step) {  // step是递归的下限
    double mid = (l + r) / 2;
    double fl = simpson(l, mid), fr = simpson(mid, r);
    if (abs(fl + fr - ans) <= 15 * eps && step < 0)
        return fl + fr + (fl + fr - ans) / 15;  // 足够相似的话就直接返回
    return asr(l, mid, eps / 2, fl, step - 1) +
           asr(mid, r, eps / 2, fr, step - 1);  // 否则分割成两段递归求解
}

double calc(double l, double r, double eps) {
    return asr(l, r, eps, simpson(l, r), 12);
}
\end{lstlisting}

\algo{多项式}
\begin{lstlisting}
struct cp {
    db x, y;
    cp(db real=0, db imag=0): x(real), y(imag) { };
    cp operator+ (cp b) const { return {x+b.x, y+b.y}; }
    cp operator- (cp b) const { return {x-b.x, y-b.y}; }
    cp operator* (cp b) const { return {x*b.x - y*b.y, x*b.y + y*b.x}; }
};
using vcp = vector<cp>;
using Poly = vector<int>;
namespace FFT {
    const db PI = acos(-1.0);
    vcp Omega(int L) {
        vcp w(L);
        w[1] = 1;
        for (int i = 2; i < L; i <<= 1) {
            auto w0 = w.begin() + i / 2;
            auto w1 = w.begin() + i;
            cp wn(cos(PI/i), sin(PI/i));
            for (int j = 0; j < i; j += 2) {
                w1[j] = w0[j>>1];
                w1[j+1] = w1[j] * wn;
            }
        }
        return w;
    }
    auto W = Omega(1 << 21);  // NOLINT
    void DIF(cp *a, int n) {
        cp x, y;
        for (int k = n>>1; k; k >>= 1) {
            for (int i = 0; i < n; i += k<<1) {
                for (int j = 0; j < k; ++j) {
                    x = a[i+j];
                    y = a[i+j+k];
                    a[i+j+k] = (a[i+j]-y) * W[k+j];
                    a[i+j] = x+y;
                }
            }
        }
    }
    void IDIT(cp *a, int n) {
        cp x, y;
        for (int k = 1; k < n; k <<= 1) {
            for (int i = 0; i < n; i += k << 1) {
                for (int j = 0; j < k; ++j) {
                    x = a[i+j];
                    y = a[i+j+k] * W[k+j];
                    a[i+j+k] = x-y;
                    a[i+j] = x+y;
                }
            }
        }
        const db Inv = 1.0 / n;
        for (int i = 0; i < n; ++i) {
            a[i].x *= Inv;
            a[i].y *= Inv;
        }
        reverse(a+1, a+n);
    }
}

namespace Polynomial {
    void DFT(vcp &a) { FFT::DIF(a.data(), a.size()); }
    void IDFT(vcp &a) { FFT::IDIT(a.data(), a.size()); }
    int norm(int n) { return 1 << (__lg(n-1) + 1); }
    vcp &dot(vcp &a, vcp &b) {
        for (int i = 0; i < int(a.size()); ++i) {
            a[i] = a[i] * b[i];
        }
        return a;
    }
    Poly operator+ (Poly a, Poly b) {
        int maxlen = max(a.size(), b.size());
        Poly ans(maxlen + 1);
        a.resize(maxlen + 1);
        b.resize(maxlen + 1);
        for (int i = 0; i < maxlen; ++i) { ans[i] = a[i] + b[i]; }
        return ans;
    }
    Poly operator* (ll k, Poly a) {
        Poly ans;
        for (auto i: a) { ans.push_back(k*i); }
        return ans;
    }
    Poly operator* (Poly a, Poly b) {
        int n = a.size() + b.size() - 1;
        vcp c(norm(n));
        for (int i = 0; i < int(a.size()); ++i) { c[i].x = a[i]; }
        for (int i = 0; i < int(b.size()); ++i) { c[i].y = b[i]; }
        DFT(c);
        dot(c, c);
        IDFT(c);
        a.resize(n);
        for (int i = 0; i < n; ++i) {
            a[i] = int(c[i].y * 0.5 + 0.5);
        }
        return a;
    }
}
using namespace Polynomial;
\end{lstlisting}

\algo{快速幂}
\begin{lstlisting}
int qmi(int m, int k, int p)
{
    int res = 1 % p, t = m;
    while (k)
    {
        if (k&1) res = res * t % p;
        t = t * t % p;
        k >>= 1;
    }
    return res;
}
\end{lstlisting}

\algo{乘法逆元}
\begin{lstlisting}
#include<bits/stdc++.h>
using namespace std;
typedef long long ll;

ll fpm(ll x, ll power, ll mod) {
    x %= mod;
    ll ans = 1;
    for (; power; power >>= 1, (x *= x) %= mod)
    	if(power & 1) (ans *= x) %= mod;
    return ans;
}
int main() {
    ios::sync_with_stdio(false);
    cin.tie(nullptr);
    int n,p;
    cin>>n>>p;
    for(int a=1;a<=n;a++){
	cout<< fpm(a, p - 2, p)<<endl;
    }//x为a在mod p意义下的逆元
}
\end{lstlisting}

\algo{组合数求法}
\begin{lstlisting}
#include <iostream>

using namespace std;

typedef long long LL;

LL qmi(LL a,int k, int p)
{
    LL res = 1;
    while(k)
    {
        if (k & 1) res = res * a % p;
        k >>= 1;
        a = a * a % p;
    }
    return res;
}

int C(int a,int b,int p)
{
    if(a < b) return 0;

    LL res = 1;
    for(int i = 1, j = a; i <= b; i ++ , j -- )
    {
        res = res * j % p;
        res = res * qmi(i, p - 2, p) % p;
    }
    return res;
}

LL lucas(LL a, LL b, int p)
{
    if(a < p && b < p) return C(a, b, p);
    return C(a % p, b % p, p) * lucas(a / p, b / p, p) % p;
}

int main()
{
    int n;
    cin >> n;
    while(n --)
    {
        LL a,b;
        int p;
        cin >> a >> b >> p;
        cout << lucas(a, b, p) << endl;
    }
    return 0;
}
\end{lstlisting}

\algo{线性基}
\begin{lstlisting}
//普通消元法
void insert(long long x)
{
    for(int i=63;~i;i--) if((x>>i)&1ll)
    {
        if(!place[i]) {place[i]=x;return;}
        else x^=place[i];
    }
}
void insert(int *num)
{
    for(int j=63;~j;j--)
    {
        int p=0;
        for(int i=1;(!p)&&i<=n;i++)
            if((!v[i])&&((num[i]>>j)&1ll)) p=i;
        if(!p) {place[j]=0;continue;}
        v[p]=true,place[j]=num[p];
        for(int i=1;i<=n;i++) if(i!=p&&((num[i]>>j)&1ll))
            num[i]^=num[p];
    }
}
bool check(long long x)
{
    for(int i=63;~i;i--)
        if((x>>i)&1ll) x^=place[i];
    return x==0;
}
struct Base
{
    long long place[N];
    void insert(long long x)
    {
        for(int i=63;~i;i--) if((x>>i)&1ll)
        {
            if(!place[i]) {place[i]=x;return;}
            else x^=place[i];
        }
    }
};
void merge(Base A,Base &B)
{
    for(int i=63;~i;i--) if(A.place[i])
        B.insert(A.place[i]);
}
long long mx()
{
    long long ans=0;
    for(int i=63;~i;i--)
        if((ans^place[i])>ans)
            ans^=place[i];
    return ans;
}
long long mn()
{
    for(int i=0;i<=63;i++)
        if(place[i]) return place[i];
    return -1;
}
\end{lstlisting}

\algo{矩阵快速幂及高斯消元}
\begin{lstlisting}
// 1.矩阵快速幂
#include <bits/stdc++.h>
using namespace std; 
#define int long long
using ll = long long;
const int MOD = 1e9 + 7;
// 矩阵结构体
struct Matrix {
    int n, m; // n: 行数, m: 列数
    vector<vector<ll>> a;
    Matrix(int _n, int _m) : n(_n), m(_m), a(_n + 1, vector<ll>(_m + 1, 0)) {}

    static Matrix identity(int n) {
        Matrix I(n, n);
        for (int i = 1; i <= n; ++i) {
            I.a[i][i] = 1;
        }
        return I;
    }
};
// 矩阵乘法: C = A * B
// 时间复杂度: O(n^3) 或 O(A.n * A.m * B.m)
Matrix operator*(const Matrix& A, const Matrix& B) {
    Matrix C(A.n, B.m);
    for (int i = 1; i <= A.n; ++i) {
        for (int j = 1; j <= B.m; ++j) {
            for (int k = 1; k <= A.m; ++k) {
                C.a[i][j] = (C.a[i][j] + A.a[i][k] * B.a[k][j]) % MOD;
            }
        }
    }
    return C;
}

// 矩阵快速幂: A^p, 时间复杂度: O(n^3 * log p) , A必须是方阵
Matrix power(Matrix A, ll p) {
    Matrix res = Matrix::identity(A.n);
    while (p) {
        if (p & 1) res= res * A;
        A = A * A;
        p >>= 1;
    }
    return res;
}
signed main(){
    ios::sync_with_stdio(false); 
    cin.tie(0), cout.tie(0); 
    int n, k; cin >> n >> k; 
    Matrix A(n, n);
    for(int i = 1; i <= n; i++){
        for(int j = 1; j <= n;j++){
            cin >> A.a[i][j]; 
        }
    }
    Matrix res = power(A, k); 
    for(int i = 1; i <= n; i++){
        for(int j = 1; j <= n; j++){
            cout << res.a[i][j] << " ";  
        }
        cout << endl; 
    }
}
\end{lstlisting}

\algo{线性筛求质因子个数（Omega函数）}
\begin{lstlisting}
const int MAXN = 1e7 + 10;
vector<int> primes;
vector<int> omega(MAXN, 0);
vector<bool> is_comp(MAXN, false);
void sieve_omega(int n) {
    omega[1] = 0;
    for (int i = 2; i <= n; ++i) {
        if (!is_comp[i]) { primes.push_back(i); omega[i] = 1; }
        for (int p : primes) {
            long long nxt = 1LL * i * p;
            if (nxt > n) break;
            is_comp[nxt] = true;
            omega[nxt] = omega[i] + 1;
            if (i % p == 0) break;
        }
    }
}
\end{lstlisting}

\algo{线性筛(O(n))}
\begin{lstlisting}
int n = 100;
vector<int> minp(n + 1, 0), prime;
//minp存的是该数的最小质因数
for(int i = 2; i <= n; i++){
 if(minp[i] == 0){
 minp[i] = i;
 prime.push_back(i);
 }
 for(auto p : prime){
    if(i * p > n) break;
    minp[i * p] = p;
    if(minp[i] == p) break;
 }
}
\end{lstlisting}

\algo{位运算的部分性质}
\begin{lstlisting}
1. 基本代数性质
交换律
c
a & b == b & a
a | b == b | a
a ^ b == b ^ a
结合律
c
(a & b) & c == a & (b & c)
(a | b) | c == a | (b | c)
(a ^ b) ^ c == a ^ (b ^ c)
幂等律
c
a & a == a
a | a == a
a ^ a == 0
注意：异或不是幂等的。

单位元、零元
运算	单位元	零元
&	M 全 1	0
|	0	M 全 1
^	0	无零元，每个数的逆元是自身
例如：

c
a & M == a
a | 0 == a
a ^ 0 == a
a ^ a == 0
互补律
c
a & ~a == 0
a | ~a == M
a ^ ~a == M
双重非
c
~~a == a
吸收律
c
a & (a | b) == a
a | (a & b) == a
分配律
c
a & (b | c) == (a & b) | (a & c)
a | (b & c) == (a | b) & (a | c)
a & (b ^ c) == (a & b) ^ (a & c)
注意：^ 对 & 一般没有分配律，即：

c
a ^ (b & c) != (a ^ b) & (a ^ c)
德摩根定律
c
~(a & b) == ~a | ~b
~(a | b) == ~a & ~b
异或取反也有常用转换：

c
~(a ^ b) == ~a ^ b == a ^ ~b
~a ^ ~b == a ^ b
2. 与、或、非、异或之间的转换
与和或互相转换
c
a | b == ~(~a & ~b)
a & b == ~(~a | ~b)
异或用与、或、非表示
c
a ^ b == (a | b) & ~(a & b)
a ^ b == (a & ~b) | (~a & b)
a ^ b == (a | b) ^ (a & b)
a ^ b == (a | b) - (a & b)
或、与用异或表示
c
a | b == a ^ b ^ (a & b)
a & b == a ^ b ^ (a | b)
非与加减法的关系
c
~a == -a - 1
-a == ~a + 1
a + ~a == M
a ^ M == ~a
3. 与算术运算的关系
加法
c
a + b == (a ^ b) + 2 * (a & b)
a + b == (a | b) + (a & b)
其中：

a ^ b 是不考虑进位的和；

a & b 是进位。

减法
c
a - b == a + ~b + 1
-a == ~a + 1
与、或、异或的算术关系
c
a | b == a + b - (a & b)
a & b == a + b - (a | b)
a ^ b == a + b - 2 * (a & b)

a & b == (a + b - (a ^ b)) / 2
a | b == (a + b + (a ^ b)) / 2
不重叠位的加法
如果：

c
a & b == 0
那么：

c
a + b == a | b == a ^ b
4. 移位与乘除的关系
左移
c
a << n == a * 2^n
在固定位宽中相当于对 2^w 取模。

右移
对于无符号数或非负数：

c
a >> n == a / 2^n   // 向下取整
对于有符号补码数，算术右移等价于向下取整：

c
a >> n == floor(a / 2^n)
例如 -3 >> 1 == -2，而不是 C/C++ 中 -3 / 2 == -1，因为 C 语言除法向零取整。

取模 2^n
c
a % 2^n == a & ((1 << n) - 1)
清零低 n 位
c
a & ~((1 << n) - 1)
循环移位
以位宽 w 为例：

c
// 循环左移 n 位
(a << n) | (a >>> (w - n))

// 循环右移 n 位
(a >>> n) | (a << (w - n))
5. 位运算与集合表示
一个整数的二进制位可以表示一个集合：第 i 位为 1 表示元素 i 在集合中。

c
交集：  a & b
并集：  a | b
对称差：a ^ b
补集：  ~a & M
差集：  a & ~b
判断子集：

c
a 是 b 的子集： (a & b) == a
            或： (a | b) == b
            或： (a & ~b) == 0
集合操作：

c
判断元素 i：   (a >> i) & 1
添加元素 i：   a | (1 << i)
删除元素 i：   a & ~(1 << i)
翻转元素 i：   a ^ (1 << i)
枚举一个集合的所有子集：

c
for (int s = a; s; s = (s - 1) & a) {
    // s 是 a 的一个非空子集
}
// 最后还要处理空集 0
6. 常用位操作恒等式
目标	表达式
判断奇偶	x & 1
乘 2	x << 1
除 2	x >> 1
判断是否为 2 的幂	x > 0 && (x & (x - 1)) == 0
获取第 k 位	(x >> k) & 1
设置第 k 位为 1	x | (1 << k)
清除第 k 位	x & ~(1 << k)
翻转第 k 位	x ^ (1 << k)
取最低位 1	x & -x
清除最低位 1	x & (x - 1)
取最低位 0	~x & (x + 1)
将最低位 0 变为 1	x | (x + 1)
将末尾连续 1 清零	x & (x + 1)
将末尾连续 0 置 1	x | (x - 1)
得到最低位 1 及以下全 1 掩码	x ^ (x - 1)
判断相等	(a ^ b) == 0
交换 a、b	a ^= b; b ^= a; a ^= b;
平均值防溢出	(a & b) + ((a ^ b) >> 1)
相反数	-x == ~x + 1
绝对值，w 位有符号数	mask = x >> (w - 1); abs = (x ^ mask) - mask;
统计 1 的个数	while (x) { x &= x - 1; count++; }
7. 注意事项
有符号与无符号
右移分算术右移和逻辑右移。算术右移补符号位，逻辑右移补 0。

移位位数
C/C++ 中移位位数超过等于位宽是未定义行为；Java 等语言会取模。

优先级
位运算优先级容易出错，建议始终加括号，例如：

c
if ((x & 1) == 0)
补码性质
补码下：

-x == ~x + 1
~x == -x - 1
逻辑运算与位运算不同
&&、||、! 是逻辑运算，结果只有 0 或 1，且会短路；
&、|、^、~ 是按位运算。
如果 x,y 只取 0/1，则：

x && y == x & y
x || y == x | y
!x == x ^ 1 == 1 - x
\end{lstlisting}

\section{字符串}

\algo{前缀函数与KMP}
\begin{lstlisting}
vector<int> prefixFunction(string s)
{   // fail : 1-index
    int n = s.length();
    vector<int> fail(n + 1);
    for (int i = 1, j = 0; i < n; ++i) {
        while (j && s[i] != s[j]) { j = fail[j]; }
        j += (s[i] == s[j]);
        fail[i + 1] = j;
    }
    return fail;
}
void kmp()
{
	string s, t;
	cin >> s >> t;
	int n = s.length();
	int m = t.length();
	auto fail = prefixFunction(t + "#" + s);
	for (int i = m + 2; i <= n + m + 1; ++i) {
	    if (fail[i] == m) {
	        cout << i - 2*m << "\n";
	    }
	}
	for (int i = 1; i <= m; ++i) { cout << fail[i] << " \n"[i == m]; }
}
\end{lstlisting}

\algo{Manacher}
\begin{lstlisting}
vector<int> manacher(string s)
{
    vector<int> d(s.length());
    d[0] = 1;
    int r=0, mid=0;
    for (int i = 1; i < s.length()-1; ++i) {  // 最后一个字符为防越界使用，无需访问
        d[i] = 1;
        if (r >= i) { d[i] = min(d[2*mid-i], r-i+1); }
        while (s[i-d[i]] == s[i+d[i]]) { ++d[i]; }  // 需要保证不会越界
        if (i + d[i] > r) { 
            r = i + d[i] - 1;
            mid = i;
        }
    }
    return d;
}
void solve()
{
    string s, t;
    cin >> s;
    for (auto c: s) {
        t.push_back('#');
        t.push_back(c);
    }
    t = "."+t+"#!";  // 首位添加两个不同的字符保证不会越界
    vector<int> d=manacher(t);
    cout << (*max_element(d.begin(), d.end()))-1 << endl;
}
\end{lstlisting}

\algo{字符串哈希}
\begin{lstlisting}
typedef pair<ll,ll> hs;
const ll MOD1=1e9+7, MOD2=1e9+9;
const hs p = { 117, 131 };
hs operator+ (hs a, hs b) {
    return hs{ (a.first+b.first)%MOD1, (a.second+b.second)%MOD2 };
}
hs operator- (hs a, hs b) {
    return hs{ (a.first-b.first+MOD1)%MOD1, (a.second-b.second+MOD2)%MOD2 };
}
hs operator* (hs a, hs b) {
    return hs{ a.first*b.first%MOD1, a.second*b.second%MOD2 };
}
struct Hash {
    int n;
    vector<hs> hs1, hs2, Pow;
    Hash(string s) {  // 0-index
        n = s.length();
        hs1.resize(n + 2);
        hs2.resize(n + 2);
        Pow.resize(n + 2);
        Pow[0] = { 1LL, 1LL };
        for (int i = 1; i <= n; ++i) {
            Pow[i] = Pow[i-1] * p;
        }
        for (int i = 1; i <= n; ++i) {
            hs1[i] = hs1[i-1] * p + hs{s[i-1]-'a'+1, s[i-1]-'a'+1};
        }
        for (int i = n; i; --i) {
            hs2[i] = hs2[i+1] * p + hs{s[i-1]-'a'+1, s[i-1]-'a'+1};
        }
    }
    hs get1(int l, int r) {
        return hs1[r] - hs1[l-1] * Pow[r-l+1];
    }
    hs get2(int l, int r) {
        return hs2[l] - hs2[r+1] * Pow[r-l+1];
    }
};
\end{lstlisting}

\algo{AC自动机}
\begin{lstlisting}
struct AhoCorasick {
    static constexpr int ALPHABET = 26;
    struct Node {
        int len;
        int link;
        array<int,ALPHABET> next;
        Node(): len{0}, link{0}, next{} { }
    };

    vector<Node> t;

    AhoCorasick() {
        init();
    }

    void init() {
        t.assign(2, Node());
        t[0].next.fill(1);
        t[0].len = -1;
    }

    int newNode() {
        t.emplace_back();
        return t.size() - 1;
    }

    int add(const string &a) {
        int p = 1;
        for (auto c: a) {
            int x = c - 'a';
            if (t[p].next[x] == 0) {
                t[p].next[x] = newNode();
                t[t[p].next[x]].len = t[p].len + 1;
            }
            p = t[p].next[x];
        }
        return p;
    }

    void work() {
        queue<int> q;
        q.push(1);

        while (!q.empty()) {
            int x = q.front();
            q.pop();

            for (int i = 0; i < ALPHABET; ++i) {
                if (t[x].next[i] == 0) {
                    t[x].next[i] = t[t[x].link].next[i];
                } else {
                    t[t[x].next[i]].link = t[t[x].link].next[i];
                    q.push(t[x].next[i]);
                }
            }
        }
    }

    int next(int p, int x) {
        return t[p].next[x];
    }

    int link(int p) {
        return t[p].link;
    }

    int len(int p) {
        return t[p].len;
    }

    int size() {
        return t.size();
    }
};

constexpr int N = 2e6+5;
vector<int> G[N];
int cnt[N];

void dfs(int u) 
{
    for (auto v: G[u]) {
        dfs(v);
        cnt[u] += cnt[v];
    }
}

void solve()
{
	int n;
	cin >> n;

    AhoCorasick AC;
    vector<int> pos(n + 1);
    for (int i = 1; i <= n; ++i) {
        string s;
        cin >> s;
        pos[i] = AC.add(s);
    }
    AC.work();

    string T;
    cin >> T;
    int p = 1;
    for (int i = 0; i < int(T.size()); ++i) {
        int x = T[i] - 'a';
        p = AC.next(p, x);
        ++cnt[p];
    }

    for (int i = 1; i < AC.size(); ++i) {
        G[AC.link(i)].push_back(i);
    }
    dfs(1);

    for (int i = 1; i <= n; ++i) {
        cout << cnt[pos[i]] << "\n";
    }
}
\end{lstlisting}

\algo{后缀自动机SAM}
\begin{lstlisting}
//edited by piaoyun from some other's code 
//必须#define int long long 

struct SAM {
    static const int N=1e6+5, S=28;
    int tot=1, last=1, link[N<<1], ch[N<<1][S], len[N<<1], endpos[N<<1];
	// 总点数 tot，点的 index 属于 [1-tot]，空串/根为 1 
	// last 为上一次插入的点
	// link 为点的 parent 树父节点 / 最长 出现位置与自己不同 的后缀
	// ch[n][s] 指节点 n 末尾加字符 s 所转移到的点
	// len 指该节点的串的 最长长度，注意到 最短长度 等于 len[link[n]] + 1 即父节点最长 + 1
	// endpos[n] 参考 get_endpos() 的注释 
    void clear() {
        for (int i = 0; i <= tot; ++i) {
            link[i] = len[i] = endpos[i] = 0;
            for (int k = 0; k < S; ++k) { ch[i][k] = 0; }
        }
        tot = 1; 
        last = 1;
    }

    // 延长一个字符，通常为 [1-26] 
    void extend(int w) {
        int p = ++tot, x = last, r, q;
        endpos[p] = 1;
        for (len[last=p]=len[x]+1; x&&!ch[x][w]; x=link[x]) { ch[x][w]=p; }
        if (!x) { link[p] = 1; }
        else if (len[x] + 1 == len[q=ch[x][w]]) { link[p]=q; }
        else {
            link[r=++tot] = link[q];
            memcpy(ch[r], ch[q], sizeof(ch[r]));
            len[r] = len[x] + 1;
            link[p] = link[q] = r;
            for (; x && ch[x][w] == q; x = link[x]) { ch[x][w] = r; }
        }
    }

    // 注意 vector 占用的空间 
    vector<int> p[N<<1];  //建立 parent 树，以便从上到下 dfs 
    void dfs(int u) {
        int v;
        for (int i = 0; i < int(p[u].size()); ++i) {
            v = p[u][i];
            dfs(v);
            endpos[u] += endpos[v];
        }
    }

    //注意！在使用该方法前，endpos[] 代表每个点作为“终结点”的次数
	//使用该方法后，endpos[] 指在串中出现总次数，即原数组的子树求和
    void get_endpos() {
        for (int i = 1; i <= tot; ++i) { p[i].clear(); }
        for (int i = 2; i <= tot; ++i) {
            p[link[i]].push_back(i);
        }
        dfs(1);
        for (int i = 1; i <= tot; ++i) { p[i].clear(); }
    }

    // 在您不确定是否有抄写错误时再使用该方法 
    // 必须在输入任何数据前自检，此前的数据会被清空 
    static const int STC = 998244353;
    void self_test() {
        clear();
        for (int i = 1; i <= 1000; ++i) { extend(i * i % 26 + 1); }
        ll tmp = 107 * last + 301 * tot;
        for (int i = 1; i <= tot; ++i) {
            tmp = (tmp * 33 + link[i] * 101 + len[i] * 97) % STC;
            for (int k = 1; k < S; ++k) { tmp = (tmp + k * ch[i][k]) % STC; }
        }
        assert("stage 1" && tmp == 393281314);   // stage1 : 检查建树是否正确
        tmp = 0;
        get_endpos();
        for (int i = 1; i <= tot; ++i) { tmp = (tmp * 33 + endpos[i]) % STC; }
        assert("stage 2" && tmp == 178417668);   // stage2 : 检查endpos计算是否正确，如果您修改了endpos[]的含义则会报错
        cout << "Self Test Passed. Remember to delete this function's use. \n";
        clear(); 
    }

	// 调试时可调用 
    void debug_print() {
        for (int i = 1; i <= tot; ++i) {
            cerr << "node : " << i << " father : " << link[i] << " endpos : " << endpos[i] << " len : " << len[i] << "\n"; 
        }
    }

    ll solve() {
		// 在这里输入你自己的解题逻辑
        ll ans = 0;
        get_endpos();
        for (int i = 1; i <= tot; ++i) {
            if (endpos[i] >= 2) { ans = max(ans, (ll)endpos[i]*len[i]); }
        }
        return ans;
    }
} sam;

void solve()
{
    // sam.self_test();
    string s; cin >> s;
    for (auto c: s) { sam.extend(c-'a'+1); }
    cout << sam.solve() << "\n";
}
\end{lstlisting}

\algo{KMP算法}
\begin{lstlisting}
#include<bits/stdc++.h>
using namespace std;
void get_nxt(string s, int nxt[]){
    int j = 0;
    nxt[0] = 0;
    for (int i = 1; i < s.length(); i++)
    {
        while(j>0&&s[i]!=s[j]) j=nxt[j-1];
        if(s[i]==s[j])j++;
        nxt[i]=j;
    }
}
void solve(){
    int cnt=0;
    string s, t;
    cin >> s >> t;
    int* nxt = new int[t.length()];
    get_nxt(t, nxt);
    int j=0;
    for (int i = 0; i < s.length(); i++) //这里的s代表主串
    {
        while(j > 0 && s[i] != t[j]) j = nxt[j-1];
        if(s[i]==t[j]) j++;
        if(j==t.length()){
            cout<<i-t.length()+2<<endl;
        }
    }
    for (int i = 0; i < t.length(); i++){
        cout<<nxt[i]<<' ';
    }
}
int main(){
    solve();
}
\end{lstlisting}

\algo{Trie树}
\begin{lstlisting}
#include<bits/stdc++.h>
using namespace std;
int tree[3000005][70],pass[3000005],endp[3000005],cnt=1;
int getidx(char ch){
// 映射函数
if(ch>='a'&&ch<='z') return ch-'a';
if(ch>='A'&&ch<='Z') return ch-'A'+26;
return ch-'0'+52;    
}
void insert(string s){//插入操作
 int cur=1;
 pass[cur]++;
    for(int i=0;i<s.length();i++){
        int idx=getidx(s[i]);
        if(tree[cur][idx]==0) {
            tree[cur][idx]=++cnt;
        }
        cur=tree[cur][idx];
        pass[cur]++;
    }
    endp[cur]++;
}
int search_whole_s(string s){
    //查询整个字符串出现次数
    int cur=1;
    for(int i=0;i<s.length();i++){
        int idx=getidx(s[i]);
        if(tree[cur][idx]==0) return 0;
        cur=tree[cur][idx];
    }
    return endp[cur];
}
int search_prefix_s(string s){
    //查询s作为前缀字符串出现次数
    int cur=1;
    for(int i=0;i<s.length();i++){
        int idx=getidx(s[i]);
        if(tree[cur][idx]==0) return 0;
        cur=tree[cur][idx];
    }
    return pass[cur];
}
void delete1(string s)
{
    if(search_whole_s(s)>0)
    {
    int cur=1;
    for(int i=0;i<s.length();i++)
    {
    int idx=getidx(s[i]);
    if(--pass[tree[cur][idx]]==0)
    {
    tree[cur][idx]==0;
    return;
    }
    cur=tree[cur][idx];
    }
    endp[cur]--;
}
}
void clear()
{
    for(int i=1;i<=cnt;i++)
    {
    for(int j=0;j<=69;j++) tree[i][j]=0;
    pass[i]=0;
    endp[i]=0;
    }
    cnt=1;
}
\end{lstlisting}

\algo{string常用函数}
\begin{lstlisting}
三、基本属性
s.size();      // 长度（常用）
s.length();    // 同 size()
s.empty();     // 是否为空
四、访问与修改
s[i];          // 访问（不检查越界）
s.at(i);       // 访问（会检查越界）

s.front();     // 第一个字符
s.back();      // 最后一个字符
五、拼接与修改
1. 拼接
s += "abc";
s = s + "abc";
2. 插入
s.insert(2, "abc");   // 在下标2插入
3. 删除
s.erase(2, 3);        // 从下标2开始删除3个字符
4. 替换
s.replace(2, 3, "abc");   // 替换区间
六、子串操作
string t = s.substr(pos, len);  // 从pos开始，长度len

示例：

string s = "abcdef";
s.substr(2, 3);   // "cde"
七、查找操作（重点）
1. 查找子串
s.find("abc");     // 返回第一次出现的位置

找不到返回：

string::npos

示例：

if (s.find("abc") != string::npos) {
    // 找到了
}
2. 反向查找
s.rfind("abc");    // 从后往前找
3. 查找字符
s.find('a');
八、比较
if (s == t) {}
if (s < t) {}   // 字典序比较

底层是字典序（lexicographical order）

九、排序
#include <algorithm>

sort(s.begin(), s.end());

示例：

string s = "dcba";
sort(s.begin(), s.end());  // "abcd"
十、反转
reverse(s.begin(), s.end());
十一、大小写转换
#include <cctype>

for (char &c : s) c = toupper(c);
for (char &c : s) c = tolower(c);
十二、统计字符
int cnt = count(s.begin(), s.end(), 'a');
\end{lstlisting}

\section{动态规划}

\algo{数位DP}
\begin{lstlisting}
#include <bits/stdc++.h>
using namespace std;
using ll = long long;
using db = long double;
using i128 = __int128;

const int L=18;
int a[L+5];
ll f[L+5][15];
ll dfs(int p, bool limit, bool lead0, int last)
{
	if (!p) { return 1; }
	auto& now = f[p][last];
	if (!limit && !lead0 && ~now) {
		return now;
	}
	int up = limit ? a[p] : 9;
	ll res = 0;
	for (int i = 0; i <= up; ++i) {
		if (!lead0 && abs(last-i) < 2) { continue; }
		res += dfs(p-1, limit&&(i==up), lead0&&(i==0), i);
	}
	if (!limit && !lead0) {
		now = res;
	}
	return res;
}
ll query(ll x)
{
	int len = 0;
	for (; x; x /= 10) {
		a[++len] = x % 10;
	}
	return dfs(len, true, true, 0);
}

void solve()
{
	ll a, b;
	cin >> a >> b;
	cout << query(b) - query(a-1) << "\n";
}

int main()
{
	ios::sync_with_stdio(0); cin.tie(0); cout.tie(0); 
	memset(f, -1, sizeof(f));  // 不要忘记设为 -1
	int t = 1;
	// cin >> t;
	while (t--) { solve(); }
}
\end{lstlisting}

\algo{单调队列优化（滑动窗口）}
\begin{lstlisting}
auto work = [&](auto&& cmp) -> void {
	deque<pair<int,int>> q;

	auto update = [&](int id, int x) -> void {
		while (!q.empty() && cmp(x, q.back().second)) {
			q.pop_back();
		}
		q.emplace_back(id, x);
	};

	for (int i = 1; i <= k; ++i) {
		update(i, a[i]);
	}
	for (int i = 1, j = i+k-1; j <= n; ++i, ++j) {
		while (!q.empty() && q.front().first < i) { q.pop_front(); }
		update(j, a[j]);
		cout << q.front().second << " \n"[j == n];
	}
};
\end{lstlisting}

\algo{无向图上回路计数}
\begin{lstlisting}
// CF11D
for (int i = 1; i <= n; ++i) {
	f[1<<(i-1)][i] = 1;
}
for (int s = 0; s < (1<<n); ++s) {
	for (int i = 1; i <= n; ++i) {
		if (!f[s][i]) { continue; }
		for (int j = 1; j <= n; ++j) {
			if (!G[i][j]) { continue; }
			if ((s&-s) > 1<<(j-1)) { continue; }  // 起点不能改
			if(s>>(j-1)&1) {
				if ((s&-s) == 1<<(j-1)) {  // 走回起点
					ans += f[s][i];
				}
			} else {
				f[s|(1<<(j-1))][j] += f[s][i];
			}
		}
	}
}
cout << (ans-m)/2 << "\n";  // 去除无向边，环正走反走算一个
\end{lstlisting}

\algo{SOSDP}
\begin{lstlisting}
for(int k=0;k<=19;k++)
        for(int i=0;i<=maxn;i++)
            if((i>>k)&1) g[i] += g[i^(1<<k)];   // i 是超集，有第 k 位
for(int k=0;k<=19;k++)
        for(int i=0;i<=maxn;i++)
            if(!((i>>k)&1)) g[i] += g[i^(1<<k)];   // i 是子集，无第 k 位
\end{lstlisting}

\algo{树形DP}
\begin{lstlisting}
#include<bits/stdc++.h>
//#define int long long
#define in(a) a = read_int()
using namespace std;
const int N = 2e5 + 5;
const int inf = INT_MAX;
inline int read_int() {
	int x = 0;
	char ch = getchar();
	bool f = 0;
	while('9' < ch || ch < '0') f |= ch == '-' , ch = getchar();
	while('0' <= ch && ch <= '9') x = (x << 3) + (x << 1) + ch - '0' , ch = getchar();
	return f ? -x : x;
}
struct Edge {
	int n , v , w;
} edge[N << 1];
int head[N] , eid;
void eadd(int u , int v , int w) {
	edge[++ eid].n = head[u] , edge[eid].v = v , edge[head[u] = eid].w = w;
}
int n;
int l[N] , r[N] , ans[N];
void dfs(int u , int fa) { // 树形 DP 部分 
	l[u] = 1; // 初始化左边界 
	for(int i = head[u] ; i ; i = edge[i].n) {
		int v = edge[i].v , w = edge[i].w;
		if(v != fa) {
			r[v] = w - 1; // 初始化右边界 
			dfs(v , u);
			r[u] = min(r[u] , w - l[v]); // 转移过程 
			l[u] = max(l[u] , w - r[v]);
		}
	}
}
queue<int>q;
signed main() {
	in(n);
	for(int i = 1 ; i < n ; i ++) {
		int u , v , w;
		in(u) , in(v) , in(w);
		eadd(u , v , w) , eadd(v , u , w);
	}
	r[1] = 1e9; // 初始化根节点的右边界 
	dfs(1 , 0); // 求每个点权的取值范围 
	for(int i = 1 ; i <= eid ; i += 2) { // 判无解部分 
		int u = edge[i].v , v = edge[i + 1].v , w = edge[i].w;
		if(r[u] + r[v] < w || l[u] > r[u] || l[v] > r[v]) { 
            cout<<"NO";
			return 0;
		}
	}
	ans[1] = l[1]; // 若有解，根节点点权取左边界 
	q.push(1);
	while(!q.empty()) { // 广搜部分 
		int u = q.front();
		q.pop();
		for(int i = head[u] ; i ; i = edge[i].n) {
			int v = edge[i].v , w = edge[i].w;
			if(ans[v]) continue;
			ans[v] = w - ans[u]; // 递推求解
		}
	}
	cout<<"YES\n"; // 打印答案 
	for(int i = 1 ; i <= n ; i ++) cout<<ans[i]<<' ';
	return 0;
    
}
/*
#include<bits/stdc++.h>
using namespace std;
const int maxn=10000;
int cnt=1,n,head[maxn],yes[maxn],no[maxn],num[maxn],b[maxn];
struct graph{
    int next,to;
}g[maxn];
void add(int u,int v){
    g[cnt].next=head[u];
    g[cnt].to=v;
    head[u]=cnt++;
}
void dfs(int i){
    no[i]=0;
    yes[i]=num[i];
    for(int ei=head[i],v;ei;ei=g[ei].next){
        v=g[ei].to;
        dfs(v);
        no[i]+=max(yes[v],no[v]);
        yes[i]+=no[v];
    }
}
int main(){
    cin>>n;
    for(int i=1;i<=n;i++) cin>>num[i];
    for(int i=1;i<n;i++){
        int u,v;
        cin>>u>>v;
        b[u]=1;
        add(v,u);
        
    }
    for(int i=1;i<=n;i++){
        if(!b[i]){
            dfs(i);
            cout<<max(yes[i],no[i])<<endl;
            break;
        }
    }
}
*/
\end{lstlisting}

\algo{最长不下降子序列}
\begin{lstlisting}
#include <bits/stdc++.h>
using namespace std;
using ll = long long;
const int N = 1e5+9;
int n,k,m,a[N],b[N],dp[N],t[N<<2];

void updata(int i,int l,int r,int n,int val){
    t[i]=max(val,t[i]);
    if(l==r) return;
  int mid=(l+r)>>1;
  if(n<=mid) updata(i<<1,l,mid,n,val);
  else updata(i<<1|1,mid+1,r,n,val);
  return;
}
int getmax(int i,int l,int r,int nl,int nr){
  if(l==nl&&r==nr){
    return t[i];
  }
  int mid=(l+r)>>1;
  if(mid>=nr) return getmax(i<<1,l,mid,nl,nr);
  else if(mid<nl) return getmax(i<<1|1,mid+1,r,nl,nr);
  else return max(getmax(i<<1,l,mid,nl,mid),getmax(i<<1|1,mid+1,r,mid+1,nr));
}
void solve(){
    cin>>n>>k;
    for(int i=1;i<=n;i++)cin>>a[i],b[i]=a[i];
    // 离散化
    sort(b+1,b+n+1);
    m=unique(b+1,b+n+1)-b-1;
    for(int i=1;i<=n;i++)a[i]=lower_bound(b+1,b+m+1,a[i])-b;
    for(int i=1;i<=n;i++){
        dp[i]=getmax(1,1,m,1,a[i])+1;
        updata(1,1,m,a[i],dp[i]);
    }
    //到这里dp 中存储的是以a[i]结尾的最长不下降子序列
    int ans=0;
    memset(t,0,sizeof t);
    for(int i=n;i>k;i--){
        ans=max(ans,dp[i-k]+k-1+getmax(1,1,m,a[i-k],m)+1); //通过再次利用线段树找到往后最长（题目要求）
        int tem=getmax(1,1,m,a[i],m)+1;
        updata(1,1,m,a[i],tem);
    }
    ans=max(ans,k+getmax(1,1,m,a[k+1],m));
    cout<<ans<<'\n';
}

// 1 2 3 ... i-k i-k+1 i-k+2 i-k+3 ... i ...

int main(){
    ios::sync_with_stdio(false),cin.tie(nullptr),cout.tie(nullptr);
    int q=1;
    //cin>>q;
    while(q--){
        solve();
    }
    return 0;
}
\end{lstlisting}

\section{计算几何}

\algo{define}
\begin{lstlisting}
using ll = long long;
using db = double;
using T = ll;   // 或者 double

// 如果 T 是整数
int sgn(ll x) {
    return (x > 0) - (x < 0);
}

// 如果 T 是浮点数，用这个版本
/*
const double eps = 1e-9;
int sgn(double x) {
    if (fabs(x) < eps) return 0;
    return x < 0 ? -1 : 1;
}
*/
\end{lstlisting}

\algo{向量}
\begin{lstlisting}
using T = ll;
struct Point {
	T x, y;
    Point(T x=0, T y=0): x(x), y(y) {}
	friend istream& operator>> (istream& is, Point& p) { return is>>p.x>>p.y; }
	friend ostream& operator<< (ostream& os, Point& p) { return os<<p.x<<" "<<p.y; }
	db ang() { return atan2(y, x); } 
	bool operator== (const Point &p) const { return sgn(x-p.x)==0 && sgn(y-p.y)==0; }
	bool operator< (const Point &p) const { return x<p.x || (x==p.x && y<p.y); }
	Point operator+ (const Point &p) const { return Point(x+p.x, y+p.y); }
    Point operator- (const Point &p) const { return Point(x-p.x, y-p.y); }
	Point operator* (const T &k) const { return Point(k*x, k*y); }
	Point operator/ (const T &k) const { return Point(x/k, y/k); }
	T operator* (const Point &p) const { return x*p.x + y*p.y; }
	T operator^ (const Point &p) const { return x*p.y - y*p.x; }  // 叉乘时打括号
	T len2() { return (*this)*(*this); }
	db len() { return sqrtl(len2()); }  // 等价于 hypotl(x, y)
	Point rot(db ang) { return Point(x*cos(ang)-y*sin(ang), x*sin(ang)+y*cos(ang)); }
	Point trunc(db l) { return (*this) * (l/len()); }
};
using Vector = Point;

// 判断 c 是否在 ab 的逆时针方向
int toLeft(Point a, Point b, Point c) { return sgn((b-a)^(c-a)); }

// 二维向量夹角
db getAngle(Vector a, Vector b) { return fabs(atan2(fabs(a^b), a*b)); }
\end{lstlisting}

\algo{点线操作}
\begin{lstlisting}
struct Line {
    Point a, b;
    Vector v;
	Line(Point a, Point b): a(a), b(b) { v = b - a; }
	int toLeft(Point p) { return sgn(v^(p-a)); }
	bool onSegment(Point p) { return (toLeft(p)==0) && ((p-a)*(p-b)<=0); }
	db distToLine(Point p) { return fabs((v^(p-a))/v.len()); }
	Point proj(Point p) { return a+v*((v*(p-a))/(v*v)); }
	// 用整数操作判断直线与点 C 的距离是否小于（等于）某个值 r
	// 常见于判断直线与圆的位置关系
	// 注意叉乘再相乘可能会爆 ll，所以这里用 __int128
	bool disLess(Point p, i128 r) {
	    i128 xmult = (v^(p-a));
	    return xmult*xmult <= r*r*v.len2();
	}
};

// 判断两条直线是否平行
bool isParallel(Line l1, Line l2)
{
	return sgn(l1.v^l2.v)==0 && sgn(l1.v^(l1.a-l2.a))!=0;
}

// 判断两条线段是否相交
// 先做两次跨立实验判断是否规范相交，再特判至少三点共线的情况
bool hasIntersection(Line L1, Line L2)
{
	// 特判三点共线的情况
    if (L1.onSegment(L2.a)) { return true; }
    if (L1.onSegment(L2.b)) { return true; }
    if (L2.onSegment(L1.a)) { return true; }
    if (L2.onSegment(L1.b)) { return true; }
    // 跨立实验判断两直线是否规范相交
    // 线段端点分布在直线两侧，“线段与线段/线段与直线/直线与直线”是否相交对应修改即可
	function<bool(Line,Point,Point)> crossStanding = [](Line L, Point a, Point b){
        return L.toLeft(a) * L.toLeft(b) == -1;
    };
    return crossStanding(L1, L2.a, L2.b) && crossStanding(L2, L1.a, L1.b);
}

// 求两直线交点
// 求之前应该先判断有没有交点，避免精度误差，同时还要特判共线的情况
Point getIntersection(Line L1, Line L2)
{
    Point P1=L1.a, P2=L2.a;
    Point v1=L1.v, v2=L2.v;
    return P1 + v1 * ( (v2^(P1-P2)) / (v1^v2) );
}

//点P到线段AB的距离公式
double distToSeg(Point p, Point a, Point b) {
    if(a == b) return (p-a).len();
    Point v1=b-a, v2=p-a, v3=p-b;
    if(sgn(v1*v2) < 0) return v2.len();
    if(sgn(v1*v3) > 0) return v3.len();
    return Line(a,b).distToLine(p);
}
\end{lstlisting}

\algo{多边形与凸包}
\begin{lstlisting}
struct Polygon: vector<Point> {
    using vector<Point>::vector;  // 直接使用 vector 的构造函数
    size_t nxt(size_t i) { return i+1==size()? 0: i+1; }
    size_t pre(size_t i) { return i==0? size()-1: i-1; }
	// 求多边形面积
	// 为避免精度误差，可以返回 2*S，即去掉 "0.5*()"
	db area() {
        db res = 0;
        for (size_t i = 0; i < size(); ++i) {
            res += 0.5 * ((*this)[i] ^ (*this)[nxt(i)]);
        }
        return fabs(res);
    }
    // Sunday's algorithm 光线回转法
    // 在整数范围内判断一个点是否在多边形（不保证凸包）内部
    // 亦可用于求某点相对于该多边形的回转数
    pair<bool,int> winding(Point p) {
        int wn = 0, n = this->size();
        for (int i = 0; i < n; ++i) {
            if (Line((*this)[i], (*this)[(i+1)%n]).onSegment(p)) {
            	return { true, -inf };
            }
            int k = Line((*this)[i], (*this)[(i+1)%n]).toLeft(p);
            int d1 = sgn((*this)[i].y - p.y);
            int d2 = sgn((*this)[(i + 1) % n].y - p.y);
            if (k>0 && d1<=0 && d2>0) { wn++; }
            if (k<0 && d2<=0 && d1>0) { wn--; }
        }
        return { wn!=0, wn };
    }
};
\end{lstlisting}

\algo{计算几何模板}
\begin{lstlisting}
#include <bits/stdc++.h>
using namespace std;

using point_t=long double;  //全局数据类型，可修改为 long long 等

constexpr point_t eps=1e-8;
constexpr long double PI=3.1415926535897932384l;

// 点与向量
template<typename T> struct point
{
    T x,y;

    bool operator==(const point &a) const {return (abs(x-a.x)<=eps && abs(y-a.y)<=eps);}
    bool operator<(const point &a) const {if (abs(x-a.x)<=eps) return y<a.y-eps; return x<a.x-eps;}
    bool operator>(const point &a) const {return !(*this<a || *this==a);}
    point operator+(const point &a) const {return {x+a.x,y+a.y};}
    point operator-(const point &a) const {return {x-a.x,y-a.y};}
    point operator-() const {return {-x,-y};}
    point operator*(const T k) const {return {k*x,k*y};}
    point operator/(const T k) const {return {x/k,y/k};}
    T operator*(const point &a) const {return x*a.x+y*a.y;}  // 点积
    T operator^(const point &a) const {return x*a.y-y*a.x;}  // 叉积，注意优先级
    int toleft(const point &a) const {const auto t=(*this)^a; return (t>eps)-(t<-eps);}  // to-left 测试
    T len2() const {return (*this)*(*this);}  // 向量长度的平方
    T dis2(const point &a) const {return (a-(*this)).len2();}  // 两点距离的平方

    // 涉及浮点数
    long double len() const {return sqrtl(len2());}  // 向量长度
    long double dis(const point &a) const {return sqrtl(dis2(a));}  // 两点距离
    long double ang(const point &a) const {return acosl(max(-1.0l,min(1.0l,((*this)*a)/(len()*a.len()))));}  // 向量夹角
    point rot(const long double rad) const {return {x*cos(rad)-y*sin(rad),x*sin(rad)+y*cos(rad)};}  // 逆时针旋转（给定角度）
    point rot(const long double cosr,const long double sinr) const {return {x*cosr-y*sinr,x*sinr+y*cosr};}  // 逆时针旋转（给定角度的正弦与余弦）
};

using Point=point<point_t>;

// 极角排序
struct argcmp
{
    bool operator()(const Point &a,const Point &b) const
    {
        const auto quad=[](const Point &a)
        {
            if (a.y<-eps) return 1;
            if (a.y>eps) return 4;
            if (a.x<-eps) return 5;
            if (a.x>eps) return 3;
            return 2;
        };
        const int qa=quad(a),qb=quad(b);
        if (qa!=qb) return qa<qb;
        const auto t=a^b;
        return t>eps;
    }
};

// 直线
template<typename T> struct line
{
    point<T> p,v;  // p 为直线上一点，v 为方向向量

    bool operator==(const line &a) const {return v.toleft(a.v)==0 && v.toleft(p-a.p)==0;}
    int toleft(const point<T> &a) const {return v.toleft(a-p);}  // to-left 测试
    bool operator<(const line &a) const  // 半平面交算法定义的排序
    {
        if (abs(v^a.v)<=eps && v*a.v>=-eps) return toleft(a.p)==-1;
        return argcmp()(v,a.v);
    }

    // 涉及浮点数
    point<T> inter(const line &a) const {return p+v*((a.v^(p-a.p))/(v^a.v));}  // 直线交点
    long double dis(const point<T> &a) const {return abs(v^(a-p))/v.len();}  // 点到直线距离
    point<T> proj(const point<T> &a) const {return p+v*((v*(a-p))/(v*v));}  // 点在直线上的投影
};

using Line=line<point_t>;

//线段
template<typename T> struct segment
{
    point<T> a,b;

    bool operator<(const segment &s) const {return make_pair(a,b)<make_pair(s.a,s.b);}

    // 判定性函数建议在整数域使用

    // 判断点是否在线段上
    // -1 点在线段端点 | 0 点不在线段上 | 1 点严格在线段上
    int is_on(const point<T> &p) const  
    {
        if (p==a || p==b) return -1;
        return (p-a).toleft(p-b)==0 && (p-a)*(p-b)<-eps;
    }

    // 判断线段直线是否相交
    // -1 直线经过线段端点 | 0 线段和直线不相交 | 1 线段和直线严格相交
    int is_inter(const line<T> &l) const
    {
        if (l.toleft(a)==0 || l.toleft(b)==0) return -1;
        return l.toleft(a)!=l.toleft(b);
    }
    
    // 判断两线段是否相交
    // -1 在某一线段端点处相交 | 0 两线段不相交 | 1 两线段严格相交
    int is_inter(const segment<T> &s) const
    {
        if (is_on(s.a) || is_on(s.b) || s.is_on(a) || s.is_on(b)) return -1;
        const line<T> l{a,b-a},ls{s.a,s.b-s.a};
        return l.toleft(s.a)*l.toleft(s.b)==-1 && ls.toleft(a)*ls.toleft(b)==-1;
    }

    // 点到线段距离
    long double dis(const point<T> &p) const
    {
        if ((p-a)*(b-a)<-eps || (p-b)*(a-b)<-eps) return min(p.dis(a),p.dis(b));
        const line<T> l{a,b-a};
        return l.dis(p);
    }

    // 两线段间距离
    long double dis(const segment<T> &s) const
    {
        if (is_inter(s)) return 0;
        return min({dis(s.a),dis(s.b),s.dis(a),s.dis(b)});
    }
};

using Segment=segment<point_t>;

// 多边形
template<typename T> struct polygon
{
    vector<point<T>> p;  // 以逆时针顺序存储

    size_t nxt(const size_t i) const {return i==p.size()-1?0:i+1;}
    size_t pre(const size_t i) const {return i==0?p.size()-1:i-1;}
    
    // 回转数
    // 返回值第一项表示点是否在多边形边上
    // 对于狭义多边形，回转数为 0 表示点在多边形外，否则点在多边形内
    pair<bool,int> winding(const point<T> &a) const
    {
        int cnt=0;
        for (size_t i=0;i<p.size();i++)
        {
            const point<T> u=p[i],v=p[nxt(i)];
            if (abs((a-u)^(a-v))<=eps && (a-u)*(a-v)<=eps) return {true,0};
            if (abs(u.y-v.y)<=eps) continue;
            const Line uv={u,v-u};
            if (u.y<v.y-eps && uv.toleft(a)<=0) continue;
            if (u.y>v.y+eps && uv.toleft(a)>=0) continue;
            if (u.y<a.y-eps && v.y>=a.y-eps) cnt++;
            if (u.y>=a.y-eps && v.y<a.y-eps) cnt--;
        }
        return {false,cnt};
    }

    // 多边形面积的两倍
    // 可用于判断点的存储顺序是顺时针或逆时针
    T area() const
    {
        T sum=0;
        for (size_t i=0;i<p.size();i++) sum+=p[i]^p[nxt(i)];
        return sum;
    }

    // 多边形的周长
    long double circ() const
    {
        long double sum=0;
        for (size_t i=0;i<p.size();i++) sum+=p[i].dis(p[nxt(i)]);
        return sum;
    }
};

using Polygon=polygon<point_t>;

//凸多边形
template<typename T> struct convex: polygon<T>
{
    // 闵可夫斯基和
    convex operator+(const convex &c) const  
    {
        const auto &p=this->p;
        vector<Segment> e1(p.size()),e2(c.p.size()),edge(p.size()+c.p.size());
        vector<point<T>> res; res.reserve(p.size()+c.p.size());
        const auto cmp=[](const Segment &u,const Segment &v) {return argcmp()(u.b-u.a,v.b-v.a);};
        for (size_t i=0;i<p.size();i++) e1[i]={p[i],p[this->nxt(i)]};
        for (size_t i=0;i<c.p.size();i++) e2[i]={c.p[i],c.p[c.nxt(i)]};
        rotate(e1.begin(),min_element(e1.begin(),e1.end(),cmp),e1.end());
        rotate(e2.begin(),min_element(e2.begin(),e2.end(),cmp),e2.end());
        merge(e1.begin(),e1.end(),e2.begin(),e2.end(),edge.begin(),cmp);
        const auto check=[](const vector<point<T>> &res,const point<T> &u)
        {
            const auto back1=res.back(),back2=*prev(res.end(),2);
            return (back1-back2).toleft(u-back1)==0 && (back1-back2)*(u-back1)>=-eps;
        };
        auto u=e1[0].a+e2[0].a;
        for (const auto &v:edge)
        {
            while (res.size()>1 && check(res,u)) res.pop_back();
            res.push_back(u);
            u=u+v.b-v.a;
        }
        if (res.size()>1 && check(res,res[0])) res.pop_back();
        return {res};
    }

    // 旋转卡壳
    // func 为更新答案的函数，可以根据题目调整位置
    template<typename F> void rotcaliper(const F &func) const
    {
        const auto &p=this->p;
        const auto area=[](const point<T> &u,const point<T> &v,const point<T> &w){return (w-u)^(w-v);};
        for (size_t i=0,j=1;i<p.size();i++)
        {
            const auto nxti=this->nxt(i);
            func(p[i],p[nxti],p[j]);
            while (area(p[this->nxt(j)],p[i],p[nxti])>=area(p[j],p[i],p[nxti]))
            {
                j=this->nxt(j);
                func(p[i],p[nxti],p[j]);
            }
        }
    }

    // 凸多边形的直径的平方
    T diameter2() const
    {
        const auto &p=this->p;
        if (p.size()==1) return 0;
        if (p.size()==2) return p[0].dis2(p[1]);
        T ans=0;
        auto func=[&](const point<T> &u,const point<T> &v,const point<T> &w){ans=max({ans,w.dis2(u),w.dis2(v)});};
        rotcaliper(func);
        return ans;
    }
    
    // 判断点是否在凸多边形内
    // 复杂度 O(logn)
    // -1 点在多边形边上 | 0 点在多边形外 | 1 点在多边形内
    int is_in(const point<T> &a) const
    {
        const auto &p=this->p;
        if (p.size()==1) return a==p[0]?-1:0;
        if (p.size()==2) return segment<T>{p[0],p[1]}.is_on(a)?-1:0; 
        if (a==p[0]) return -1;
        if ((p[1]-p[0]).toleft(a-p[0])==-1 || (p.back()-p[0]).toleft(a-p[0])==1) return 0;
        const auto cmp=[&](const Point &u,const Point &v){return (u-p[0]).toleft(v-p[0])==1;};
        const size_t i=lower_bound(p.begin()+1,p.end(),a,cmp)-p.begin();
        if (i==1) return segment<T>{p[0],p[i]}.is_on(a)?-1:0;
        if (i==p.size()-1 && segment<T>{p[0],p[i]}.is_on(a)) return -1;
        if (segment<T>{p[i-1],p[i]}.is_on(a)) return -1;
        return (p[i]-p[i-1]).toleft(a-p[i-1])>0;
    }

    // 凸多边形关于某一方向的极点
    // 复杂度 O(logn)
    // 参考资料：https://codeforces.com/blog/entry/48868
    template<typename F> size_t extreme(const F &dir) const
    {
        const auto &p=this->p;
        const auto check=[&](const size_t i){return dir(p[i]).toleft(p[this->nxt(i)]-p[i])>=0;};
        const auto dir0=dir(p[0]); const auto check0=check(0);
        if (!check0 && check(p.size()-1)) return 0;
        const auto cmp=[&](const Point &v)
        {
            const size_t vi=&v-p.data();
            if (vi==0) return 1;
            const auto checkv=check(vi);
            const auto t=dir0.toleft(v-p[0]);
            if (vi==1 && checkv==check0 && t==0) return 1;
            return checkv^(checkv==check0 && t<=0);
        };
        return partition_point(p.begin(),p.end(),cmp)-p.begin();
    }

    // 过凸多边形外一点求凸多边形的切线，返回切点下标
    // 复杂度 O(logn)
    // 必须保证点在多边形外
    pair<size_t,size_t> tangent(const point<T> &a) const
    {
        const size_t i=extreme([&](const point<T> &u){return u-a;});
        const size_t j=extreme([&](const point<T> &u){return a-u;});
        return {i,j};
    }

    // 求平行于给定直线的凸多边形的切线，返回切点下标
    // 复杂度 O(logn)
    pair<size_t,size_t> tangent(const line<T> &a) const
    {
        const size_t i=extreme([&](...){return a.v;});
        const size_t j=extreme([&](...){return -a.v;});
        return {i,j};
    }
};

using Convex=convex<point_t>;

// 圆
struct Circle
{
    Point c;
    long double r;

    bool operator==(const Circle &a) const {return c==a.c && abs(r-a.r)<=eps;}
    long double circ() const {return 2*PI*r;}  // 周长
    long double area() const {return PI*r*r;}  // 面积

    // 点与圆的关系
    // -1 圆上 | 0 圆外 | 1 圆内
    int is_in(const Point &p) const {const long double d=p.dis(c); return abs(d-r)<=eps?-1:d<r-eps;}

    // 直线与圆关系
    // 0 相离 | 1 相切 | 2 相交
    int relation(const Line &l) const
    {
        const long double d=l.dis(c);
        if (d>r+eps) return 0;
        if (abs(d-r)<=eps) return 1;
        return 2;
    }

    // 圆与圆关系
    // -1 相同 | 0 相离 | 1 外切 | 2 相交 | 3 内切 | 4 内含
    int relation(const Circle &a) const
    {
        if (*this==a) return -1;
        const long double d=c.dis(a.c);
        if (d>r+a.r+eps) return 0;
        if (abs(d-r-a.r)<=eps) return 1;
        if (abs(d-abs(r-a.r))<=eps) return 3;
        if (d<abs(r-a.r)-eps) return 4;
        return 2;
    }

    // 直线与圆的交点
    vector<Point> inter(const Line &l) const
    {
        const long double d=l.dis(c);
        const Point p=l.proj(c);
        const int t=relation(l);
        if (t==0) return vector<Point>();
        if (t==1) return vector<Point>{p};
        const long double k=sqrt(r*r-d*d);
        return vector<Point>{p-(l.v/l.v.len())*k,p+(l.v/l.v.len())*k};
    }

    // 圆与圆交点
    vector<Point> inter(const Circle &a) const
    {
        const long double d=c.dis(a.c);
        const int t=relation(a);
        if (t==-1 || t==0 || t==4) return vector<Point>();
        Point e=a.c-c; e=e/e.len()*r;
        if (t==1 || t==3) 
        {
            if (r*r+d*d-a.r*a.r>=-eps) return vector<Point>{c+e};
            return vector<Point>{c-e};
        }
        const long double costh=(r*r+d*d-a.r*a.r)/(2*r*d),sinth=sqrt(1-costh*costh);
        return vector<Point>{c+e.rot(costh,-sinth),c+e.rot(costh,sinth)};
    }

    // 圆与圆交面积
    long double inter_area(const Circle &a) const
    {
        const long double d=c.dis(a.c);
        const int t=relation(a);
        if (t==-1) return area();
        if (t<2) return 0;
        if (t>2) return min(area(),a.area());
        const long double costh1=(r*r+d*d-a.r*a.r)/(2*r*d),costh2=(a.r*a.r+d*d-r*r)/(2*a.r*d);
        const long double sinth1=sqrt(1-costh1*costh1),sinth2=sqrt(1-costh2*costh2);
        const long double th1=acos(costh1),th2=acos(costh2);
        return r*r*(th1-costh1*sinth1)+a.r*a.r*(th2-costh2*sinth2);
    }

    // 过圆外一点圆的切线
    vector<Line> tangent(const Point &a) const
    {
        const int t=is_in(a);
        if (t==1) return vector<Line>();
        if (t==-1)
        {
            const Point v={-(a-c).y,(a-c).x};
            return vector<Line>{{a,v}};
        }
        Point e=a-c; e=e/e.len()*r;
        const long double costh=r/c.dis(a),sinth=sqrt(1-costh*costh);
        const Point t1=c+e.rot(costh,-sinth),t2=c+e.rot(costh,sinth);
        return vector<Line>{{a,t1-a},{a,t2-a}};
    }

    // 两圆的公切线
    vector<Line> tangent(const Circle &a) const
    {
        const int t=relation(a);
        vector<Line> lines;
        if (t==-1 || t==4) return lines;
        if (t==1 || t==3)
        {
            const Point p=inter(a)[0],v={-(a.c-c).y,(a.c-c).x};
            lines.push_back({p,v});
        }
        const long double d=c.dis(a.c);
        const Point e=(a.c-c)/(a.c-c).len();
        if (t<=2)
        {
            const long double costh=(r-a.r)/d,sinth=sqrt(1-costh*costh);
            const Point d1=e.rot(costh,-sinth),d2=e.rot(costh,sinth);
            const Point u1=c+d1*r,u2=c+d2*r,v1=a.c+d1*a.r,v2=a.c+d2*a.r;
            lines.push_back({u1,v1-u1}); lines.push_back({u2,v2-u2});
        }
        if (t==0)
        {
            const long double costh=(r+a.r)/d,sinth=sqrt(1-costh*costh);
            const Point d1=e.rot(costh,-sinth),d2=e.rot(costh,sinth);
            const Point u1=c+d1*r,u2=c+d2*r,v1=a.c-d1*a.r,v2=a.c-d2*a.r;
            lines.push_back({u1,v1-u1}); lines.push_back({u2,v2-u2});
        }
        return lines;
    }

    // 圆的反演
    tuple<int,Circle,Line> inverse(const Line &l) const
    {
        const Circle null_c={{0.0,0.0},0.0};
        const Line null_l={{0.0,0.0},{0.0,0.0}};
        if (l.toleft(c)==0) return {2,null_c,l};
        const Point v=l.toleft(c)==1?Point{l.v.y,-l.v.x}:Point{-l.v.y,l.v.x};
        const long double d=r*r/l.dis(c);
        const Point p=c+v/v.len()*d;
        return {1,{(c+p)/2,d/2},null_l};
    }
    
    tuple<int,Circle,Line> inverse(const Circle &a) const
    {
        const Circle null_c={{0.0,0.0},0.0};
        const Line null_l={{0.0,0.0},{0.0,0.0}};
        const Point v=a.c-c;
        if (a.is_in(c)==-1)
        {
            const long double d=r*r/(a.r+a.r);
            const Point p=c+v/v.len()*d;
            return {2,null_c,{p,{-v.y,v.x}}};
        }
        if (c==a.c) return {1,{c,r*r/a.r},null_l};
        const long double d1=r*r/(c.dis(a.c)-a.r),d2=r*r/(c.dis(a.c)+a.r);
        const Point p=c+v/v.len()*d1,q=c+v/v.len()*d2;
        return {1,{(p+q)/2,p.dis(q)/2},null_l};
    }
};

// 圆与多边形面积交
long double area_inter(const Circle &circ,const Polygon &poly)
{
    const auto cal=[](const Circle &circ,const Point &a,const Point &b)
    {
        if ((a-circ.c).toleft(b-circ.c)==0) return 0.0l;
        const auto ina=circ.is_in(a),inb=circ.is_in(b);
        const Line ab={a,b-a};
        if (ina && inb) return ((a-circ.c)^(b-circ.c))/2;
        if (ina && !inb)
        {
            const auto t=circ.inter(ab);
            const Point p=t.size()==1?t[0]:t[1];
            const long double ans=((a-circ.c)^(p-circ.c))/2;
            const long double th=(p-circ.c).ang(b-circ.c);
            const long double d=circ.r*circ.r*th/2;
            if ((a-circ.c).toleft(b-circ.c)==1) return ans+d;
            return ans-d;
        }
        if (!ina && inb)
        {
            const Point p=circ.inter(ab)[0];
            const long double ans=((p-circ.c)^(b-circ.c))/2;
            const long double th=(a-circ.c).ang(p-circ.c);
            const long double d=circ.r*circ.r*th/2;
            if ((a-circ.c).toleft(b-circ.c)==1) return ans+d;
            return ans-d;
        }
        const auto p=circ.inter(ab);
        if (p.size()==2 && Segment{a,b}.dis(circ.c)<=circ.r+eps)
        {
            const long double ans=((p[0]-circ.c)^(p[1]-circ.c))/2;
            const long double th1=(a-circ.c).ang(p[0]-circ.c),th2=(b-circ.c).ang(p[1]-circ.c);
            const long double d1=circ.r*circ.r*th1/2,d2=circ.r*circ.r*th2/2;
            if ((a-circ.c).toleft(b-circ.c)==1) return ans+d1+d2;
            return ans-d1-d2;
        }
        const long double th=(a-circ.c).ang(b-circ.c);
        if ((a-circ.c).toleft(b-circ.c)==1) return circ.r*circ.r*th/2;
        return -circ.r*circ.r*th/2;
    };

    long double ans=0;
    for (size_t i=0;i<poly.p.size();i++)
    {
        const Point a=poly.p[i],b=poly.p[poly.nxt(i)];
        ans+=cal(circ,a,b);
    }
    return ans;
}

// 点集的凸包
// Andrew 算法，复杂度 O(nlogn)
Convex convexhull(vector<Point> p)
{
    vector<Point> st;
    if (p.empty()) return Convex{st};
    sort(p.begin(),p.end());
    const auto check=[](const vector<Point> &st,const Point &u)
    {
        const auto back1=st.back(),back2=*prev(st.end(),2);
        return (back1-back2).toleft(u-back1)<=0;
    };
    for (const Point &u:p)
    {
        while (st.size()>1 && check(st,u)) st.pop_back();
        st.push_back(u);
    }
    size_t k=st.size();
    p.pop_back(); reverse(p.begin(),p.end());
    for (const Point &u:p)
    {
        while (st.size()>k && check(st,u)) st.pop_back();
        st.push_back(u);
    }
    st.pop_back();
    return Convex{st};
}

// 半平面交
// 排序增量法，复杂度 O(nlogn)
// 输入与返回值都是用直线表示的半平面集合
vector<Line> halfinter(vector<Line> l, const point_t lim=1e9)
{
    const auto check=[](const Line &a,const Line &b,const Line &c){return a.toleft(b.inter(c))<0;};
    l.push_back({{-lim,0},{0,-1}}); l.push_back({{0,-lim},{1,0}});
    l.push_back({{lim,0},{0,1}}); l.push_back({{0,lim},{-1,0}});
    sort(l.begin(),l.end());
    deque<Line> q;
    for (size_t i=0;i<l.size();i++)
    {
        if (i>0 && l[i-1].v.toleft(l[i].v)==0 && l[i-1].v*l[i].v>eps) continue;
        while (q.size()>1 && check(l[i],q.back(),q[q.size()-2])) q.pop_back();
        while (q.size()>1 && check(l[i],q[0],q[1])) q.pop_front();
        if (!q.empty() && q.back().v.toleft(l[i].v)<=0) return vector<Line>();
        q.push_back(l[i]);
    }
    while (q.size()>1 && check(q[0],q.back(),q[q.size()-2])) q.pop_back();
    while (q.size()>1 && check(q.back(),q[0],q[1])) q.pop_front();
    return vector<Line>(q.begin(),q.end());
}

// 点集形成的最小最大三角形
// 极角序扫描线，复杂度 O(n^2logn)
// 最大三角形问题可以使用凸包与旋转卡壳做到 O(n^2)
pair<point_t,point_t> minmax_triangle(const vector<Point> &vec)
{
    if (vec.size()<=2) return {0,0};
    vector<pair<int,int>> evt;
    evt.reserve(vec.size()*vec.size());
    point_t maxans=0,minans=numeric_limits<point_t>::max();
    for (size_t i=0;i<vec.size();i++)
    {
        for (size_t j=0;j<vec.size();j++)
        {
            if (i==j) continue;
            if (vec[i]==vec[j]) minans=0;
            else evt.push_back({i,j});
        }
    }
    sort(evt.begin(),evt.end(),[&](const pair<int,int> &u,const pair<int,int> &v)
    {
        const Point du=vec[u.second]-vec[u.first],dv=vec[v.second]-vec[v.first];
        return argcmp()({du.y,-du.x},{dv.y,-dv.x});
    });
    vector<size_t> vx(vec.size()),pos(vec.size());
    for (size_t i=0;i<vec.size();i++) vx[i]=i;
    sort(vx.begin(),vx.end(),[&](int x,int y){return vec[x]<vec[y];});
    for (size_t i=0;i<vx.size();i++) pos[vx[i]]=i;
    for (auto [u,v]:evt)
    {
        const size_t i=pos[u],j=pos[v];
        const size_t l=min(i,j),r=max(i,j);
        const Point vecu=vec[u],vecv=vec[v];
        if (l>0) minans=min(minans,abs((vec[vx[l-1]]-vecu)^(vec[vx[l-1]]-vecv)));
        if (r<vx.size()-1) minans=min(minans,abs((vec[vx[r+1]]-vecu)^(vec[vx[r+1]]-vecv)));
        maxans=max({maxans,abs((vec[vx[0]]-vecu)^(vec[vx[0]]-vecv)),abs((vec[vx.back()]-vecu)^(vec[vx.back()]-vecv))});
        if (i<j) swap(vx[i],vx[j]),pos[u]=j,pos[v]=i;
    }
    return {minans,maxans};
}

// 判断多条线段是否有交点
// 扫描线，复杂度 O(nlogn)
bool segs_inter(const vector<Segment> &segs)
{
    if (segs.empty()) return false;
    using seq_t=tuple<point_t,int,Segment>;
    const auto seqcmp=[](const seq_t &u, const seq_t &v)
    {
        const auto [u0,u1,u2]=u;
        const auto [v0,v1,v2]=v;
        if (abs(u0-v0)<=eps) return make_pair(u1,u2)<make_pair(v1,v2);
        return u0<v0-eps;
    };
    vector<seq_t> seq;
    for (auto seg:segs)
    {
        if (seg.a.x>seg.b.x+eps) swap(seg.a,seg.b);
        seq.push_back({seg.a.x,0,seg});
        seq.push_back({seg.b.x,1,seg});
    }
    sort(seq.begin(),seq.end(),seqcmp);
    point_t x_now;
    auto cmp=[&](const Segment &u, const Segment &v)
    {
        if (abs(u.a.x-u.b.x)<=eps || abs(v.a.x-v.b.x)<=eps) return u.a.y<v.a.y-eps;
        return ((x_now-u.a.x)*(u.b.y-u.a.y)+u.a.y*(u.b.x-u.a.x))*(v.b.x-v.a.x)<((x_now-v.a.x)*(v.b.y-v.a.y)+v.a.y*(v.b.x-v.a.x))*(u.b.x-u.a.x)-eps;
    };
    multiset<Segment,decltype(cmp)> s{cmp};
    for (const auto [x,o,seg]:seq)
    {
        x_now=x;
        const auto it=s.lower_bound(seg);
        if (o==0)
        {
            if (it!=s.end() && seg.is_inter(*it)) return true;
            if (it!=s.begin() && seg.is_inter(*prev(it))) return true;
            s.insert(seg);
        }
        else
        {
            if (next(it)!=s.end() && it!=s.begin() && (*prev(it)).is_inter(*next(it))) return true;
            s.erase(it);
        }
    }
    return false;
}

// 多边形面积并
// 轮廓积分，复杂度 O(n^2logn)，n为边数
// ans[i] 表示被至少覆盖了 i+1 次的区域的面积
vector<long double> area_union(const vector<Polygon> &polys)
{
    const size_t siz=polys.size();
    vector<vector<pair<Point,Point>>> segs(siz);
    const auto check=[](const Point &u,const Segment &e){return !((u<e.a && u<e.b) || (u>e.a && u>e.b));};

    auto cut_edge=[&](const Segment &e,const size_t i)
    {
        const Line le{e.a,e.b-e.a};
        vector<pair<Point,int>> evt;
        evt.push_back({e.a,0}); evt.push_back({e.b,0});
        for (size_t j=0;j<polys.size();j++)
        {
            if (i==j) continue;
            const auto &pj=polys[j];
            for (size_t k=0;k<pj.p.size();k++)
            {
                const Segment s={pj.p[k],pj.p[pj.nxt(k)]};
                if (le.toleft(s.a)==0 && le.toleft(s.b)==0)
                {
                    evt.push_back({s.a,0});
                    evt.push_back({s.b,0});
                }
                else if (s.is_inter(le))
                {
                    const Line ls{s.a,s.b-s.a};
                    const Point u=le.inter(ls);
                    if (le.toleft(s.a)<0 && le.toleft(s.b)>=0) evt.push_back({u,-1});
                    else if (le.toleft(s.a)>=0 && le.toleft(s.b)<0) evt.push_back({u,1});
                }
            }
        }
        sort(evt.begin(),evt.end());
        if (e.a>e.b) reverse(evt.begin(),evt.end());
        int sum=0;
        for (size_t i=0;i<evt.size();i++)
        {
            sum+=evt[i].second;
            const Point u=evt[i].first,v=evt[i+1].first;
            if (!(u==v) && check(u,e) && check(v,e)) segs[sum].push_back({u,v});
            if (v==e.b) break;
        }
    };
    
    for (size_t i=0;i<polys.size();i++)
    {
        const auto &pi=polys[i];
        for (size_t k=0;k<pi.p.size();k++)
        {
            const Segment ei={pi.p[k],pi.p[pi.nxt(k)]};
            cut_edge(ei,i);
        }
    }
    vector<long double> ans(siz);
    for (size_t i=0;i<siz;i++)
    {
        long double sum=0;
        sort(segs[i].begin(),segs[i].end());
        int cnt=0;
        for (size_t j=0;j<segs[i].size();j++)
        {
            if (j>0 && segs[i][j]==segs[i][j-1]) segs[i+(++cnt)].push_back(segs[i][j]);
            else cnt=0,sum+=segs[i][j].first^segs[i][j].second;
        }
        ans[i]=sum/2;
    }
    return ans;
}

// 圆面积并
// 轮廓积分，复杂度 O(n^2logn)
// ans[i] 表示被至少覆盖了 i+1 次的区域的面积
vector<long double> area_union(const vector<Circle> &circs)
{
    const size_t siz=circs.size();
    using arc_t=tuple<Point,long double,long double,long double>;
    vector<vector<arc_t>> arcs(siz);
    const auto eq=[](const arc_t &u,const arc_t &v)
    {
        const auto [u1,u2,u3,u4]=u;
        const auto [v1,v2,v3,v4]=v;
        return u1==v1 && abs(u2-v2)<=eps && abs(u3-v3)<=eps && abs(u4-v4)<=eps;
    };

    auto cut_circ=[&](const Circle &ci,const size_t i)
    {
        vector<pair<long double,int>> evt;
        evt.push_back({-PI,0}); evt.push_back({PI,0});
        int init=0;
        for (size_t j=0;j<circs.size();j++)
        {
            if (i==j) continue;
            const Circle &cj=circs[j];
            if (ci.r<cj.r-eps && ci.relation(cj)>=3) init++;
            const auto inters=ci.inter(cj);
            if (inters.size()==1) evt.push_back({atan2l((inters[0]-ci.c).y,(inters[0]-ci.c).x),0});
            if (inters.size()==2)
            {
                const Point dl=inters[0]-ci.c,dr=inters[1]-ci.c;
                long double argl=atan2l(dl.y,dl.x),argr=atan2l(dr.y,dr.x);
                if (abs(argl+PI)<=eps) argl=PI;
                if (abs(argr+PI)<=eps) argr=PI;
                if (argl>argr+eps)
                {
                    evt.push_back({argl,1}); evt.push_back({PI,-1});
                    evt.push_back({-PI,1}); evt.push_back({argr,-1});
                }
                else
                {
                    evt.push_back({argl,1});
                    evt.push_back({argr,-1});
                }
            }
        }
        sort(evt.begin(),evt.end());
        int sum=init;
        for (size_t i=0;i<evt.size();i++)
        {
            sum+=evt[i].second;
            if (abs(evt[i].first-evt[i+1].first)>eps) arcs[sum].push_back({ci.c,ci.r,evt[i].first,evt[i+1].first});
            if (abs(evt[i+1].first-PI)<=eps) break;
        }
    };

    const auto oint=[](const arc_t &arc)
    {
        const auto [cc,cr,l,r]=arc;
        if (abs(r-l-PI-PI)<=eps) return 2.0l*PI*cr*cr;
        return cr*cr*(r-l)+cc.x*cr*(sin(r)-sin(l))-cc.y*cr*(cos(r)-cos(l));
    };

    for (size_t i=0;i<circs.size();i++)
    {
        const auto &ci=circs[i];
        cut_circ(ci,i);
    }
    vector<long double> ans(siz);
    for (size_t i=0;i<siz;i++)
    {
        long double sum=0;
        sort(arcs[i].begin(),arcs[i].end());
        int cnt=0;
        for (size_t j=0;j<arcs[i].size();j++)
        {
            if (j>0 && eq(arcs[i][j],arcs[i][j-1])) arcs[i+(++cnt)].push_back(arcs[i][j]);
            else cnt=0,sum+=oint(arcs[i][j]);
        }
        ans[i]=sum/2;
    }
    return ans;
}
\end{lstlisting}

\section{其它}

\algo{矩阵加速}
\begin{lstlisting}
//1939
// #include <iostream>
// #include <cstdio>
// #include <cstring>

// using namespace std;

// typedef long long LL;
// const int Mod = 1e9+7;
// int T, n;
// struct mat{
// 	LL m[5][5];
// }Ans, base;
// inline void init() {
// 	memset(Ans.m, 0, sizeof(Ans.m));
// 	for(int i=1; i<=3; i++) Ans.m[i][i] = 1;
// 	memset(base.m, 0, sizeof(base.m));
// 	base.m[1][1] = base.m[1][3] = base.m[2][1] = base.m[3][2] = 1;
// }
// inline mat mul(mat a, mat b) {
// 	mat res;
// 	memset(res.m, 0, sizeof(res.m));
// 	for(int i=1; i<=3; i++) {
// 		for(int j=1; j<=3; j++) {
// 			for(int k=1; k<=3; k++) {
// 				res.m[i][j] += (a.m[i][k] % Mod) * (b.m[k][j] % Mod);
// 				res.m[i][j] %= Mod;
// 			}
// 		}
// 	}
// 	return res;
// }
// inline void Qmat_pow(int p) {
// 	while (p) {
// 		if(p & 1) Ans = mul(Ans, base);
// 		base = mul(base, base);
// 		p >>= 1;
// 	}
// }

// int main() {
// 	scanf("%d", &T);
// 	while (T--) {
// 		scanf("%d", &n);
// 		if(n <= 3) {
// 			printf("1\n");
// 			continue;
// 		}
// 		init();
// 		Qmat_pow(n);
// 		printf("%lld\n", Ans.m[2][1]);
// 	}
// }
/*
f[i-1]       f[i]
f[i-2] ----> f[i-1]
f[i-3]       f[i-2]
f[i] = f[i-1] * 1 + f[i-2] * 0 + f[i-3] * 1
f[i-1] = f[i-1] * 1 + f[i-2] * 0 + f[i-3] * 0
f[i-2] = f[i-1] * 0 + f[i-2] * 1 + f[i-3] * 0
so
1 0 1
1 0 0
0 1 0
*/
//P6569
#include<iostream>
#include<cstdio>
#include<cstring>
#include<algorithm>
using namespace std;
struct matrix{//用结构体存矩阵敲好使哒~
	long long qwq[105][105];
    int x,y;//x,y就是矩阵的行数和列数
	matrix(){x=0,y=0;memset(qwq,0,sizeof(qwq));};//构造函数，刚听说这玩意，用着玩玩
}; 
matrix operator * (const matrix &a,const matrix &b){//重载“异或版矩阵乘法”的运算符
	matrix c;
	c.x=a.x,c.y=b.y;
	for(int i=1;i<=a.x;i++){  
		for(int j=1;j<=b.y;j++){
			for(int k=1;k<=a.y;k++)
			c.qwq[i][j]^=a.qwq[i][k]*b.qwq[k][j];//和矩阵乘法别无二致，只不过就是改成异或
		}
	}
	return c;                   
}
matrix wyx[40];//存储e矩阵的2^i次方，取这个变量名是因为想借助神仙的力量AC本题
long long f[105];//每个城市的初始魔法值（提醒：十年OI一场空……）
int main(){
	int n,m,q;
	scanf("%d%d%d",&n,&m,&q);
	for(int i=1;i<=n;i++)
	scanf("%lld",&f[i]);
	wyx[0].x=n,wyx[0].y=n;//初始化一开始的邻接矩阵（也就是2的0次方）的大小
	for(int i=1;i<=m;i++){
		int a,b;
		scanf("%d%d",&a,&b);
		wyx[0].qwq[a][b]=wyx[0].qwq[b][a]=1;//给邻接矩阵连上边
	}
	for(int i=1;i<32;i++)wyx[i]=wyx[i-1]*wyx[i-1];//利用快速幂的思想，处理出邻接矩阵2的i次方
	while(q--){
		long long x;
		scanf("%lld",&x);
        matrix ans;
		for(int i=1;i<=n;i++)ans.qwq[1][i]=f[i];//一开始就是初始f值
		ans.x=1,ans.y=n;
		for(int i=0;i<32;i++){//开始进行二进制拆分
			if((x>>i)&1)//如果这一位是1
			ans=ans*wyx[i];//就乘上对应的2的i次方
		}
		printf("%lld\n",ans.qwq[1][1]);
	}
	return 0;
}
\end{lstlisting}

\algo{并查集}
\begin{lstlisting}
//推导部分和
#include <bits/stdc++.h>
using namespace std;
typedef long long ll;
const ll N=1e5+10;
struct tree{
  ll n,w,to;
}e[N<<1];
ll head[N<<1],tot,fa[N],vis[N],sum[N];
void add(int u,int v,ll w){
  e[tot].to=v;
  e[tot].n=head[u];
  e[tot].w=w;
  head[u]=tot++;
}
int find(int i){
  if(fa[i]!=i) return fa[i]=find(fa[i]);
  return i;
}
void unite(int x,int y){
  x=find(x);
  y=find(y);
  if(x!=y) fa[x]=y;
}
void bfs(int x){
  vis[x]=1;
  queue<int>q; q.push(x);
  sum[x]=0;
  while(!q.empty()){
    int u=q.front();
    q.pop();
    for(int i=head[u];i!=-1;i=e[i].n){
      int v=e[i].to;
      if(vis[v]) continue;
      vis[v]=1;
      sum[v]=sum[u]+e[i].w;
      q.push(v);
    }
  }
}
int main()
{
  memset(head,-1,sizeof(head));
  int n,m,q;
  cin>>n>>m>>q;
  for(int i=1;i<=n;i++) fa[i]=i;
  for(int i=1;i<=m;i++){
    ll l,r,w;
    cin>>l>>r>>w;
    add(l-1,r,w);
    add(r,l-1,-w);
    unite(l-1,r);
  }
  for(int i=0;i<=n;i++){
    if(!vis[i]) bfs(i);
  }
  for(int i=1;i<=m;i++){
    int l,r;
    cin>>l>>r;
    if(find(l-1)!=find(r)) cout<<"UNKNOWN"<<endl;
    else cout<<sum[r]-sum[l-1]<<endl;
  }
  return 0;
}
\end{lstlisting}

\algo{FFT（快速傅里叶变换）}
\begin{lstlisting}
// 当vector用就可以了
#include <bits/stdc++.h>
#define fp(i, a, b) for (int i = (a), i##_ = (b) + 1; i < i##_; ++i)
#define fd(i, a, b) for (int i = (a), i##_ = (b) - 1; i > i##_; --i)
using namespace std;
using ll = int64_t;
using db = double;
/*---------------------------------------------------------------------------*/
struct cp {
    db x, y;
    cp(db real = 0, db imag = 0) : x(real), y(imag) {};
    cp operator+(cp b) const { return {x + b.x, y + b.y}; }
    cp operator-(cp b) const { return {x - b.x, y - b.y}; }
    cp operator*(cp b) const { return {x * b.x - y * b.y, x * b.y + y * b.x}; }
};
using vcp = vector<cp>;
using Poly = vector<int>;
namespace FFT {
const db pi = acos(-1);
vcp Omega(int L) {
    vcp w(L);
    w[1] = 1;
    for (int i = 2; i < L; i <<= 1) {
        auto w0 = w.begin() + i / 2, w1 = w.begin() + i;
        cp wn(cos(pi / i), sin(pi / i));
        for (int j = 0; j < i; j += 2)
            w1[j] = w0[j >> 1], w1[j + 1] = w1[j] * wn;
    }
    return w;
}
auto W = Omega(1 << 21);  // NOLINT
void DIF(cp *a, int n) {
    cp x, y;
    for (int k = n >> 1; k; k >>= 1)
        for (int i = 0; i < n; i += k << 1)
            for (int j = 0; j < k; ++j)
                x = a[i + j], y = a[i + j + k],
                a[i + j + k] = (a[i + j] - y) * W[k + j], a[i + j] = x + y;
}
void IDIT(cp *a, int n) {
    cp x, y;
    for (int k = 1; k < n; k <<= 1)
        for (int i = 0; i < n; i += k << 1)
            for (int j = 0; j < k; ++j)
                x = a[i + j], y = a[i + j + k] * W[k + j], a[i + j + k] = x - y,
                a[i + j] = x + y;
    const db Inv = 1. / n;
    fp(i, 0, n - 1) a[i].x *= Inv, a[i].y *= Inv;
    reverse(a + 1, a + n);
}
}  // namespace FFT

namespace Polynomial {
// basic operator
void DFT(vcp &a) { FFT::DIF(a.data(), a.size()); }
void IDFT(vcp &a) { FFT::IDIT(a.data(), a.size()); }
int norm(int n) { return 1 << (__lg(n - 1) + 1); }

// Poly mul
vcp &dot(vcp &a, vcp &b) {
    fp(i, 0, a.size() - 1) a[i] = a[i] * b[i];
    return a;
}
Poly operator+(Poly a, Poly b) {
    int maxlen = max(a.size(), b.size());
    Poly ans(maxlen + 1);
    a.resize(maxlen + 1), b.resize(maxlen + 1);
    for (int i = 0; i < maxlen; i++) ans[i] = a[i] + b[i];
    return ans;
}
Poly operator*(ll k, Poly a) {
    Poly ans;
    for (auto i : a) ans.push_back(k * i);
    return ans;
}
Poly operator*(Poly a, Poly b) {
    int n = a.size() + b.size() - 1;
    vcp c(norm(n));
    fp(i, 0, a.size() - 1) c[i].x = a[i];
    fp(i, 0, b.size() - 1) c[i].y = b[i];
    DFT(c), dot(c, c), IDIT(c), a.resize(n);
    fp(i, 0, n - 1) a[i] = int(c[i].y * .5 + .5);
    return a;
}
}  // namespace Polynomial
/*---------------------------------------------------------------------------*/
using namespace Polynomial;
\end{lstlisting}

\algo{NTT（数论变换）}
\begin{lstlisting}
#include <bits/stdc++.h>
#define fp(i, a, b) for (int i = (a), i##_ = (b) + 1; i < i##_; ++i)
#define fd(i, a, b) for (int i = (a), i##_ = (b) - 1; i > i##_; --i)
using namespace std;
const int maxn = 2e5 + 5, P = 998244353;
using ll = int64_t;
#define ADD(a, b) (((a) += (b)) >= P ? (a) -= P : 0)
#define SUB(a, b) (((a) -= (b)) < 0 ? (a) += P : 0)
#define MUL(a, b) (ll(a) * (b) % P)
int POW(ll a, int b = P - 2, ll x = 1) {
    for (; b; b >>= 1, a = a * a % P)
        if (b & 1) x = x * a % P;
    return x;
}

namespace NTT {
const int g = 3;
vector<int> Omega(int L) {
    int wn = POW(g, P / L);
    vector<int> w(L);
    w[L >> 1] = 1;
    fp(i, L / 2 + 1, L - 1) w[i] = MUL(w[i - 1], wn);
    fd(i, L / 2 - 1, 1) w[i] = w[i << 1];
    return w;
}
auto W = Omega(1 << 21);
void DIF(int *a, int n) {
    for (int k = n >> 1; k; k >>= 1) {
        for (int i = 0, y; i < n; i += k << 1)
            fp(j, 0, k - 1) y = a[i + j + k],
                            a[i + j + k] = MUL(a[i + j] - y + P, W[k + j]),
                            ADD(a[i + j], y);
    }
}
void IDIT(int *a, int n) {
    for (int k = 1; k < n; k <<= 1) {
        for (int i = 0, x, y; i < n; i += k << 1)
            fp(j, 0, k - 1) x = a[i + j], y = MUL(a[i + j + k], W[k + j]),
                            a[i + j + k] = x - y < 0 ? x - y + P : x - y,
                            ADD(a[i + j], y);
    }
    int Inv = P - (P - 1) / n;
    fp(i, 0, n - 1) a[i] = MUL(a[i], Inv);
    reverse(a + 1, a + n);
}
}  // namespace NTT
namespace Polynomial {
using Poly = std::vector<int>;
Poly &operator*=(Poly &a, int b) {
    for (auto &x : a) x = MUL(x, b);
    return a;
}
Poly operator*(Poly a, int b) { return a *= b; }
Poly operator*(int a, Poly b) { return b * a; }
void DFT(Poly &a) { NTT::DIF(a.data(), a.size()); }
void IDFT(Poly &a) { NTT::IDIT(a.data(), a.size()); }
int norm(int n) { return 1 << (32 - __builtin_clz(n - 1)); }
void norm(Poly &a) {
    if (!a.empty()) a.resize(norm(a.size()), 0);
}
Poly &dot(Poly &a, Poly &b) {
    fp(i, 0, a.size() - 1) a[i] = MUL(a[i], b[i]);
    return a;
}
Poly operator*(Poly a, Poly b) {
    int n = a.size() + b.size() - 1, L = norm(n);
    if (a.size() <= 8 || b.size() <= 8) {
        Poly c(n);
        fp(i, 0, a.size() - 1) fp(j, 0, b.size() - 1) c[i + j] =
            (c[i + j] + (ll)a[i] * b[j]) % P;
        return c;
    }
    a.resize(L), b.resize(L);
    DFT(a), DFT(b), dot(a, b), IDFT(a);
    return a.resize(n), a;
}
}  // namespace Polynomial
using namespace Polynomial;
\end{lstlisting}

\algo{扫描线算法}
\begin{lstlisting}
#include<bits/stdc++.h>
using namespace std;
typedef long long ll;
ll arr[15010],n,height[15010];

struct node{
    ll l,r,h;
}a[15010];

struct compare{
    bool operator()(pair<ll,ll> a,pair<ll,ll> b){
        if(a.first!=b.first) return a.first<b.first;
        return a.second>b.second;
        //大根堆
    }
};
bool cmp(node x,node y){
    return x.l<y.l;
}
int main(){
    ll l,r,h,m=0;
    while(scanf("%lld %lld %lld",&l,&h,&r)!=EOF){
        a[++m].l=l;
        a[m].r=r-1;
        a[m].h=h;
        arr[++n]=l;
        arr[++n]=r;
        arr[++n]=r-1;
    }
    sort(arr+1,arr+n+1);
    ll len=unique(arr+1,arr+n+1)-arr-1;
    for(int j=1;j<=m;j++){
        a[j].l=lower_bound(arr+1,arr+len+1,a[j].l)-arr;
        a[j].r=lower_bound(arr+1,arr+len+1,a[j].r)-arr;
    }
    sort(a+1,a+m+1,cmp);
   priority_queue<pair<ll,ll>, vector<pair<ll,ll>>, compare> heap;
   for(int i=1,j=0;i<=len;i++){
    for(;j<=m&&a[j].l<=i;j++){
        heap.push(make_pair(a[j].h,a[j].r));
    }
    while(!heap.empty()&&heap.top().second<i)
    {
        heap.pop();
    }
   
   if(!heap.empty()){
    height[i]=heap.top().first;
   }
}
    ll ansn=0;
    pair<ll,ll>ans[15010];
    for(ll i=1,pre=0;i<=len;i++){
        if(pre!=height[i]){
            ans[++ansn]=make_pair(arr[i],height[i]);
        }
        pre=height[i];
    }
    for(int i=1;i<=ansn;i++){
        cout<<ans[i].first<<' '<<ans[i].second<<' ';
    }
}
\end{lstlisting}

\algo{最长上升子序列}
\begin{lstlisting}
#include<cstdio>
#include<algorithm>
const int MAXN=200001;
int a[MAXN];
int d[MAXN];
int main()
{
    int n;
    scanf("%d",&n);
    for(int i=1;i<=n;i++)
        scanf("%d",&a[i]);
    d[1]=a[1];
    int len=1;
    for(int i=2;i<=n;i++)
    {
        if(a[i]>d[len])
            d[++len]=a[i];
        else
        {
            int j=std::lower_bound(d+1,d+len+1,a[i])-d;
            d[j]=a[i]; 
        }
    }
    printf("%d\n",len);    
    return 0;
}
\end{lstlisting}

\algo{全排列输出}
\begin{lstlisting}
#include <iostream>
#include <vector>
#include <algorithm>

int main() {
    std::vector<int> arr = {1, 2, 3};  // 示例数组
    
    // 先排序（必须步骤，确保生成所有排列）
    std::sort(arr.begin(), arr.end());
    
    // 输出所有排列组合
    do {
        for (int num : arr) {
            std::cout << num << " ";
        }
        std::cout << "\n";
    } while (std::next_permutation(arr.begin(), arr.end()));
    
    return 0;
}
\end{lstlisting}

\algo{位运算用法}
\begin{lstlisting}
/*为了更好的理解状压dp，首先介绍位运算相关的知识。

'&'符号，x&y，会将两个十进制数在二进制下进行与运算(都1为1，其余为0） 然后返回其十进制下的值。例如3(11)&2(10)=2(10)。
'|'符号，x|y，会将两个十进制数在二进制下进行或运算（都0为0，其余为1） 然后返回其十进制下的值。例如3(11)|2(10)=3(11)。
'^'符号，x^y，会将两个十进制数在二进制下进行异或运算（不同为1，其余 为0）然后返回其十进制下的值。例如3(11)^2(10)=1(01)。
'~'符号，~x，按位取反。例如~101=010。
'<<'符号，左移操作，x<<2，将x在二进制下的每一位向左移动两位，最右边用0填充，x<<2相当于让x乘以4。 '>>'符号，是右移操作，x>>1相当于给x/2，去掉x二进制下的最右一位
1.判断一个数字x二进制下第i位是不是等于1。（最低第1位）

方法：if(((1<<(i−1))&x)>0) 将1左移i-1位，相当于制造了一个只有第i位 上是1，其他位上都是0的二进制数。然后与x做与运算，如果结果>0， 说明x第i位上是1，反之则是0。

2.将一个数字x二进制下第i位更改成1。

方法：x=x|(1<<(i−1)) 证明方法与1类似。

3.将一个数字x二进制下第i位更改成0。

方法：x=x&~(1<<(i−1))

4.把一个数字二进制下最靠右的第一个1去掉。

方法：x=x&(x−1)
__builtin_popcount() 二进制数中一的个数
*/
#include<iostream>
using namespace std;
struct state {
    int st[1 << 12+1];  
    int num;          
} a[15];
int m, n;
int f[15][1 << 12+1];   
const int MOD = 100000000;
void getstate(int i, int t) {
    a[i].num = 0;
    for (int j = 0; j < (1 << n); j++) {
        if ((j & (j << 1)) || (j & t)) 
            continue;
        a[i].st[++a[i].num] = j;
    }
}
void dp() {
    for (int i = 1; i <= m; i++)
        for (int j = 1; j <= a[i].num; j++)
            f[i][j] = 0;
    for (int j = 1; j <= a[1].num; j++)
        f[1][j] = 1;
    for (int i = 2; i <= m; i++) {
        for (int j = 1; j <= a[i].num; j++) {
            int s1 = a[i].st[j];
            for (int k = 1; k <= a[i-1].num; k++) {
                int s2 = a[i-1].st[k];
                if (s1 & s2) continue; 
                f[i][j] = (f[i][j] + f[i-1][k]) % MOD;
            }
        }
    }
    int ans = 0;
    for (int j = 1; j <= a[m].num; j++)
        ans = (ans + f[m][j]) % MOD;
    cout << ans << endl;
}
int main() {
    cin >> m >> n;  
    for (int i = 1; i <= m; i++) {
        int t = 0;
        for (int j = 0; j < n; j++) {
            int x;
            cin >> x;
            t = (t << 1) + (1 - x); 
        }
        getstate(i, t); 
    }
    dp();
    return 0;
}
\end{lstlisting}

\algo{Map的使用}
\begin{lstlisting}
#include<bits/stdc++.h>
using namespace std;
int main(){
	//什么是map？map是干什么用的
	map<int,int>mp; //定义了一个名叫mp的map容器
	mp[10001]=82;mp[100002]=82;
	//输出10001信息
	cout<<mp[10001]<<endl;//可能会溢出
	//安全的方法
	if(mp.find(10005)!=mp.end()){
		cout<<mp[10005]<<endl;
	}

    else cout<<"没有此人信息"<<endl;
	//储存学生姓名对应成绩
	map<string,int>stu;
	stu["xiaoming"]=85;
	stu["xiaohong"]=77;
	cout<<stu["xiaohong"]<<endl;
	//map的遍历
    //1.
	for(auto it:stu){//将容器从头到尾遍历了一遍
		//map自动按照首位字典序从小到大排序
		cout<<it.first<<' '<<it.second<<endl;
	}
	//2.迭代器类型
	map<string,int>::iterator it;//定义map<string,int>类型的迭代器it
	for(it=stu.begin();it!=stu.end();it++){
		cout<<(*it).first<<' '<<(*it).second<<endl;
	}
}
\end{lstlisting}

\algo{优先队列（堆）}
\begin{lstlisting}
//自定义类的优先队列
struct Student {
    std::string name;
    int score;
    int age;
    
    // 重载 < 运算符，用于大根堆（按score从大到小）
    bool operator<(const Student& other) const {
        // 注意：priority_queue默认用 < 比较，但实际表现是最大堆
        // 所以这里返回 true 表示优先级更低
        return score < other.score;  // 大根堆
    }
};
priority_queue<Student>q1;//大根堆
priority_queue<int>q2;//默认大根堆
priority_queue<int,vector<int>,greater<int>>q3;//小根堆
\end{lstlisting}

\algo{扫描线+线段树（矩形面积并）}
\begin{lstlisting}
#include<bits/stdc++.h>
using namespace std;
typedef long long ll;
ll arr[15010],n,height[15010];

struct node{
    ll l,r,h;
}a[15010];

struct compare{
    bool operator()(pair<ll,ll> a,pair<ll,ll> b){
        if(a.first!=b.first) return a.first<b.first;
        return a.second>b.second;
        //大根堆
    }
};
bool cmp(node x,node y){
    return x.l<y.l;
}
int main(){
    ll l,r,h,m=0;
    while(scanf("%lld %lld %lld",&l,&h,&r)!=EOF){
        a[++m].l=l;
        a[m].r=r-1;
        a[m].h=h;
        arr[++n]=l;
        arr[++n]=r;
        arr[++n]=r-1;
    }
    sort(arr+1,arr+n+1);
    ll len=unique(arr+1,arr+n+1)-arr-1;
    for(int j=1;j<=m;j++){
        a[j].l=lower_bound(arr+1,arr+len+1,a[j].l)-arr;
        a[j].r=lower_bound(arr+1,arr+len+1,a[j].r)-arr;
    }
    sort(a+1,a+m+1,cmp);
   priority_queue<pair<ll,ll>, vector<pair<ll,ll>>, compare> heap;
   for(int i=1,j=0;i<=len;i++){
    for(;j<=m&&a[j].l<=i;j++){
        heap.push(make_pair(a[j].h,a[j].r));
    }
    while(!heap.empty()&&heap.top().second<i)
    {
        heap.pop();
    }
   
   if(!heap.empty()){
    height[i]=heap.top().first;
   }
}
    ll ansn=0;
    pair<ll,ll>ans[15010];
    for(ll i=1,pre=0;i<=len;i++){
        if(pre!=height[i]){
            ans[++ansn]=make_pair(arr[i],height[i]);
        }
        pre=height[i];
    }
    for(int i=1;i<=ansn;i++){
        cout<<ans[i].first<<' '<<ans[i].second<<' ';
    }
}
\end{lstlisting}

\algo{建树参考}
\begin{lstlisting}
#include<bits/stdc++.h>
//#define int long long
#define in(a) a = read_int()
using namespace std;
const int N = 2e5 + 5;
const int inf = INT_MAX;
inline int read_int() {
	int x = 0;
	char ch = getchar();
	bool f = 0;
	while('9' < ch || ch < '0') f |= ch == '-' , ch = getchar();
	while('0' <= ch && ch <= '9') x = (x << 3) + (x << 1) + ch - '0' , ch = getchar();
	return f ? -x : x;
}
struct Edge {
	int n , v , w;
} edge[N << 1];
int head[N] , eid;
void eadd(int u , int v , int w) {
	edge[++ eid].n = head[u] , edge[eid].v = v , edge[head[u] = eid].w = w;
}
int n;
int l[N] , r[N] , ans[N];
void dfs(int u , int fa) { // 树形 DP 部分 
	l[u] = 1; // 初始化左边界 
	for(int i = head[u] ; i ; i = edge[i].n) {
		int v = edge[i].v , w = edge[i].w;
		if(v != fa) {
			r[v] = w - 1; // 初始化右边界 
			dfs(v , u);
			r[u] = min(r[u] , w - l[v]); // 转移过程 
			l[u] = max(l[u] , w - r[v]);
		}
	}
}
queue<int>q;
signed main() {
	in(n);
	for(int i = 1 ; i < n ; i ++) {
		int u , v , w;
		in(u) , in(v) , in(w);
		eadd(u , v , w) , eadd(v , u , w);
	}
	r[1] = 1e9; // 初始化根节点的右边界 
	dfs(1 , 0); // 求每个点权的取值范围 
	for(int i = 1 ; i <= eid ; i += 2) { // 判无解部分 
		int u = edge[i].v , v = edge[i + 1].v , w = edge[i].w;
		if(r[u] + r[v] < w || l[u] > r[u] || l[v] > r[v]) { 
            cout<<"NO";
			return 0;
		}
	}
	ans[1] = l[1]; // 若有解，根节点点权取左边界 
	q.push(1);
	while(!q.empty()) { // 广搜部分 
		int u = q.front();
		q.pop();
		for(int i = head[u] ; i ; i = edge[i].n) {
			int v = edge[i].v , w = edge[i].w;
			if(ans[v]) continue;
			ans[v] = w - ans[u]; // 递推求解 
			q.push(v);
		}
	}
	cout<<"YES\n"; // 打印答案 
	for(int i = 1 ; i <= n ; i ++) cout<<ans[i]<<' ';
	return 0;
}
\end{lstlisting}

\algo{Hash树}
\begin{lstlisting}
//树同构

#include <bits/stdc++.h>
 
 typedef long long LL;
 const int N = 60, MO = 998244353;
 
 struct Edge {
     int nex, v;
 }edge[N << 1]; int tp;
 
 int e[N], n, m, turn, fr[N], p[1000010], top, siz[N], h[N], g[N];
 std::vector<int> v[N];
 bool vis[1000010];
 
 inline void add(int x, int y) {
     tp++;
     edge[tp].v = y;
     edge[tp].nex = e[x];
     e[x] = tp;
     return;
 }
 
 inline bool equal(int a, int b) {
     int len = v[a].size();
     if(len != v[b].size()) return false;
     for(int i = 0; i < len; i++) {
         if(v[a][i] != v[b][i]) return false;
     }
     return true;
 }
 
 inline void getp(int n) {
     for(int i = 2; i <= n; i++) {
         if(!vis[i]) {
             p[++top] = i;
         }
         for(int j = 1; j <= top && i * p[j] <= n; j++) {
             vis[i * p[j]] = 1;
             if(i % p[j] == 0) {
                 break;
             }
         }
     }
     return;
 }
 
 void DFS_1(int x, int f) {
     siz[x] = 1;
     h[x] = 1;
     for(int i = e[x]; i; i = edge[i].nex) {
         int y = edge[i].v;
         if(y == f) continue;
         DFS_1(y, x);
         h[x] = (h[x] + 1ll * h[y] * p[siz[y]] % MO) % MO;
         siz[x] += siz[y];
     }
     return;
 }
 
 void DFS_2(int x, int f, int V) {
     g[x] = (h[x] + 1ll * V * p[n - siz[x]] % MO) % MO;
     v[turn].push_back(g[x]);
     V = (1ll * V * p[n - siz[x]] % MO + 1) % MO;
     for(int i = e[x]; i; i = edge[i].nex) {
         int y = edge[i].v;
         if(y == f) {
             continue;
         }
         DFS_2(y, x, ((LL)V + h[x] - 1 - 1ll * h[y] * p[siz[y]] % MO + MO) % MO);
     }
     return;
 }
 
 int main() {
     getp(1000009);
     scanf("%d", &m);
     for(turn = 1; turn <= m; turn++) {
         scanf("%d", &n);
         tp = 0;
         memset(e + 1, 0, n * sizeof(int));
         for(int i = 1, x; i <= n; i++) {
             scanf("%d", &x);
             if(x) {
                 add(x, i);
                 add(i, x);
             }
         }
         DFS_1(1, 0);
         DFS_2(1, 0, 0);
         std::sort(v[turn].begin(), v[turn].end());
     }
 
     for(int i = 1; i <= m; i++) {
         fr[i] = i;
     }
     for(int i = 2; i <= m; i++) {
         for(int j = 1; j < i; j++) {
             if(equal(i, j)) {
                 fr[i] = fr[j];
                 break;
             }
         }
     }
     for(int i = 1; i <= m; i++) {
         printf("%d\n", fr[i]);
     }
     return 0;
}
\end{lstlisting}

\end{document}